% !TEX root = ../main.tex
\clearpage
\chapter{부록}
\label{ch:appendix}

% Demote headings one level under Chapter 8 (sections become 8.1.x).
\renewcommand{\dissertationSubheading}[2][]{%
  \phantomsection
  \subsection{#2}%
  \if\relax\detokenize{#1}\relax\else\label{#1}\fi
}
\renewcommand{\dissertationSubsubheading}[2][]{%
  \phantomsection
  \subsubsection{#2}%
  \if\relax\detokenize{#1}\relax\else\label{#1}\fi
}

\section{평가 처리 과정과 재현성}
\label{ch:appendix-eval}

이 부록은 평가 처리 과정(파이프라인)을 기록해요. 이 처리 과정은 \ref{ch:results}장의 바탕이 되고, 부록~\ref{sec:map-registered}에 있는 등록 분석(미리 등록한 그대로 한 분석)의 바탕도 돼요.

\dissertationSubheading[sec:pipeline-stages]{처리 과정의 단계}

평가 소스 코드는 공개 저장소 \url{https://github.com/abitocodes/phd_works}의 \texttt{1-Dissertation-Works/ERC20-Allowance-PageRank-Wallet-Reputation/experiments} 폴더에 있어요. 아래의 경로는 모두 이 폴더를 기준으로 한 상대 경로예요. 평가는 다섯 단계로 이루어져요. 처음 세 단계는 아비트럼 원(Arbitrum One)에서 스캔~A의 기간인 2023년 10월 1일부터 2026년 6월 30일까지를 대상으로 돌아가요(UTC, 양 끝 날 포함). 네 번째 단계는 2026년 7월부터 9월까지를 더해요. 다섯 번째 단계는 앞의 네 단계가 빅쿼리(BigQuery)에 남긴 표를 DefiLlama의 하루 가격, Chainlink의 가격 피드(블록체인에 올라오는 가격 정보)와 함께 읽어요. 그리고 그 표에서 각 동결일(점수를 고정하는 날) 전의 열두 달을 잘라 내요.

\begin{enumerate}
\item 계획, 코호트, 그래프: 이 연구의 로그(계약이 남기는 기록)를 하나라도 뽑기 전에 분석 계획(\texttt{docs/analysis\_plan.md})과 설정 파일을 커밋해요(커밋 \texttt{43d8045}, 2026년 10월 9일 13:41~UTC). 그 기간의 로그를 프로젝트 표 하나에 옮겨 적어요(스캔~A). 주인(허락을 해 주는 쪽) 적어도 세 명이 허락한 스펜더(spender, 허락을 받아 대신 쓰는 쪽)를 나열해요(\texttt{sql/02\_spender\_candidates.sql}). 그 스펜더들의 계정 종류를 읽고(\texttt{scripts/check\_account\_code.py}), 그 가운데 스마트 계약(블록체인 위에서 저절로 실행되는 프로그램)을 계획의 코호트(묶음)로 남겨요(\texttt{scripts/build\_cohort.py}). 코호트의 에고 이벤트(코호트 계약이 한쪽 끝에 있는 이벤트)를 뽑아요(\texttt{sql/03\_ego\_events.sql}). 그리고 이 이벤트를 모아 $T_{\mathrm{obs}}$와 $t_1$에서 두 그래프의 쌍 표(주소 쌍마다 한 행씩 모은 표)를 만들어요(\texttt{sql/04\_anchor\_pairs.sql}).
\item 점검과 일치도: \ref{sec:notation}절의 풀이 프로그램(solver) 설정으로 $T_{\mathrm{obs}}$에서 모든 방법의 점수를 매겨요. 그리고 \ref{sec:proxy-formulas}절의 점검을 계산해요. 스피어먼의 $\rho$(두 줄 세우기가 얼마나 비슷한지 나타내는 수)와 켄달의 $\tau$(두 줄 세우기가 얼마나 같은 순서인지 나타내는 수)는 95\% 짝지은 부트스트랩(다시 뽑기) 구간과 함께 보고해요(다시 뽑기 400번, 시드 42). \ref{sec:statistical-tests}절의 대비도 함께 보고해요(\texttt{scripts/run\_same\_window.py}).
\item 시간 홀드아웃(나중 기간을 떼어 두고 맞혀 보는 시험) W0: $t_1$에서 모든 방법의 점수를 매겨요. 그리고 스펜더 코호트와 트레이더(거래하는 사람) 코호트를 2026년 4월부터 6월까지의 라벨(나중에 채점에 쓰는 정답 값)에 견주어 평가해요. 규칙~A의 대비도 함께 계산해요(\texttt{scripts/run\_holdout.py}).
\item 등록 기간 W1: 설정 파일의 상태를 등록됨(registered)으로 정해 두고 등록 문서(\texttt{docs/registration\_w1.md})를 커밋해요(커밋 \texttt{0d3a057}, 2026년 10월 9일 16:49~UTC). \texttt{scripts/extract\_w1.py} 스크립트는 그 커밋 전에는 돌아가지 않아요. 이 스크립트는 두 파일의 해시(파일 내용으로 계산한, 지문 같은 값)를 추출 목록(manifest)에 적어요. 그리고 스캔~B와 7월부터 9월까지의 에고 추출을 돌려요. \texttt{scripts/run\_registered.py} 스크립트는 규칙~B, 그리고 두 화살표(엣지) 종류 어느 쪽으로도 만들지 않은 결과 값에서의 비교를 평가해요.
\item 수정 분석(등록 뒤에 금액을 달러로 바꿔 다시 한 분석; \ref{sec:amount-revision}절): \texttt{scripts/usd\_inputs.py} 스크립트는 토큰 목록과 하루 가격을 만들어요. \texttt{scripts/chainlink\_check.py} 스크립트는 그 가격을 Chainlink 가격 피드와 견주어요. \texttt{scripts/run\_usd\_pipeline.py} 스크립트는 둘을 빅쿼리에 올리고, 1단계와 4단계의 에고 표에서 다섯 토큰의 이벤트를 남겨요(\texttt{sql/19}, \texttt{sql/12}). 그다음 \texttt{scripts/run\_usd12m\_pipeline.py} 스크립트가 각 동결일의 코호트를 그 동결일 전 열두 달의 허락에서 뽑아요. 그리고 그 코호트의 쌍 표, 상수, 대리 지표(대신 재는 값), 라벨을 계산해요(\texttt{sql/20}부터 \texttt{sql/23}까지, 그리고 이 스크립트가 만든 표에 돌리는 \texttt{sql/13}부터 \texttt{sql/17}까지). 또 남긴 허락 쌍과 송금 쌍이 모두 뽑아 낸 쌍 가운데 있는지 확인해요. 그다음 \texttt{CONTRACT\_REP\_VARIANT=usd12m}으로 두면, 2단계부터 4단계까지의 스크립트가 수정 분석을 계산해요. 그리고 \texttt{scripts/run\_robustness.py}, \texttt{scripts/sybil\_model.py}, \texttt{scripts/run\_benchmark.py} 스크립트가 감쇠 계수 점검, 금액 기울기 점검, 표본 정의 점검, 시빌 모형(가짜 주소를 많이 만들어 속이는 공격을 따져 본 모형), 시간 잰 실행을 더해요. \texttt{scripts/run\_supplementary.py} 스크립트는 어느 규칙에도 속하지 않는 대비를 계산해요. 이때 저장해 둔 점수와 같은 다시 뽑은 표본을 써요. 계산하는 대비는 이래요. 활동 재시작(활발한 주소에서 더 자주 다시 출발하기)을 쓰는 엔도스랭크(EndorseRank)를 받은 허락 수(허락 입차수)와 견준 대비, 그 점수를 엔도스랭크와 견준 대비, 그리고 C-PR(결합 페이지랭크, 두 그래프를 함께 걷는 점수)을 다른 모든 점수와 두 원시 차수(화살표를 그냥 센 수)와 견준 대비예요(표~\ref{tab:cpr-contrasts}).
\end{enumerate}

이 단계들의 순서는 커밋 기록, 비용 장부(ledger), W1의 추출 목록, 그리고 수정 분석의 실행 목록(manifest)에 남아 있어요. 수정 분석 요약 파일의 요약 목록(digest)에는 그 요약 파일을 계산한 코드 버전 \texttt{0a94215}가 적혀 있어요. 요약 파일은 \texttt{950a7a4}에 커밋했어요. 시간 잰 실행은 다른 일을 하지 않는 기계 한 대에서 했어요. 이 기계에는 Intel Core Ultra~9 185H 프로세서와 64~GB 메모리가 있고, Windows~11(빌드 26200), Python 3.11.9, NumPy 2.3.2, SciPy 1.16.1을 썼어요. \texttt{scripts/export\_latex.py} 스크립트가 요약 파일로 \texttt{results/tables/}의 표와 \texttt{Figures/generated/}의 그림을 만들었어요.

실험 폴더에는 다음이 들어 있어요. 질의, 스크립트, 테스트, 버전을 매긴 설정 파일(\texttt{config/contract\_reputation.yaml}), 계획과 등록 문서와 수정 문서(\texttt{docs/}), Chainlink와 견준 결과가 붙은 토큰 목록과 하루 가격(\texttt{data/0-usd-token-list-and-prices/}), 그리고 비용 장부예요. 표의 바탕이 되는 요약 파일은 \texttt{data/3-published-results-for-thesis/revised-usd-12m/}에 있고, 등록 분석의 요약 파일은 \texttt{data/3-published-results-for-thesis/}에 있어요. 두 폴더에는 저마다 설정 파일, 계획, 등록 문서의 복사본이 있어요. 또 모든 숫자에 대해 그 숫자가 든 요약 파일과 키를 알려 주는 요약 목록 \texttt{results\_digest.md}, 그리고 \texttt{SHA256SUMS.txt}가 있어요. 그래서 처리 과정을 다시 돌리지 않고도 표의 어떤 숫자든 그 요약 파일과 맞춰 볼 수 있어요. 뽑아 낸 표와 계약마다의 점수는 90\,MB보다 작은 parquet 조각으로 나누어, 공개 구글 드라이브(Google Drive) 폴더에 두었어요. 커밋해 둔 파일 \texttt{drive-manifest.json}에는 그 폴더의 링크가 있고, 모든 조각이 체크섬(파일이 바뀌지 않았는지 확인하는 값)과 함께 적혀 있어요.

인공 자료로 만든 테스트는 풀이 프로그램, 화살표 무게(가중치), 짝지은 부트스트랩, 달러(USD) 환산을 오프라인(인터넷에 연결하지 않은 상태)으로 점검해요. 실제로 추출하려면 구글 클라우드(Google Cloud) 자격 증명과 요금을 낼 프로젝트가 있어야 해요.

\dissertationSubheading[sec:bq-extraction]{빅쿼리 추출 설계}

공개 표 \texttt{bigquery-public-data.goog\_blockchain\_arbitrum\_one\_us.logs}는 기간마다 한 번씩만 읽어요. 스캔~A는 2023년 10월 1일부터 2026년 6월 30일까지의 ERC-20(토큰 공통 규칙) \texttt{Approval} 로그와 \texttt{Transfer} 로그, GMX~V2 포지션 닫기, 그리고 Aave~V3 풀 이벤트를 달마다 나눈 프로젝트 표 하나에 옮겨 적어요. 이 논문의 어떤 분석도 Aave~V3 풀 이벤트는 쓰지 않아요. 스캔~B는 2026년 7월부터 9월까지에 대해 같은 일을 해요. 그 뒤의 질의는 모두 이 프로젝트 표들을 읽어요. 그래서 체인 전체를 훑는 비용은 기간마다 한 번만 들어요. 그리고 에고 추출은 한쪽 끝에 코호트 계약이 있는 이벤트만 남겨요. 그래서 뽑아 낸 자료의 크기는 체인 전체가 아니라 코호트를 따라가요. 부록~\ref{ch:appendix-sql}에 그 질의들이 있어요.

모든 작업은 \texttt{data/bq\_ledger.json}에 질의, 매개변수, 어림한 바이트 수와 요금이 매겨진 바이트 수, 그리고 비용과 함께 적혀 있어요. 작업 56개에 요금이 매겨진 양은 $6.76$\,TiB예요. 쓴 만큼 내는 요금, 곧 TiB당 USD~6.25로 치면 USD~42.24예요. 설정 파일에 정해 둔 예산은 USD~50이에요. 스캔~A 하나에만 $3.29$\,TiB(USD~20.55)가 매겨졌고, 스캔~B에는 $0.26$\,TiB(USD~1.64)가 매겨졌어요. 등록 분석의 작업 17개에는 USD~29.71이 매겨졌어요. 수정 분석은 작업 세 개가 USD~1.31을 들여 만든 달러 에고 표를 읽어요. 그리고 수정 분석 자신의 작업 11개에는 USD~2.49가 매겨졌어요. 달러 표에 돌린 다른 작업 12개에는 USD~3.78이 매겨졌어요. 그 가운데 세 개는 \texttt{scripts/run\_bq.py} 밖에서 돌렸고, 장부에 손으로 적어 넣었어요. 그리고 탐색적(미리 정하지 않고 살펴본)이라는 꼬리표를 붙인 작업 13개에는 USD~4.95가 매겨졌어요. 이 작업들은 이 논문이 보고하지 않는 분석과 사후 점검(결과를 본 뒤에 한 점검)에 썼어요.

\dissertationSubheading[sec:account-type-classification]{계정 종류 나누기}

코호트에 들어갈 수 있는 주소는 모두 \ref{sec:data-sources}절의 규칙에 따라 그 코드로 종류를 나누었어요. \texttt{scripts/check\_account\_code.py}는 주소마다 한 블록에서 \texttt{eth\_getCode}를 불러요. 그리고 그 블록을, 읽은 시각과 함께 적어 둬요. 표~\ref{tab:account-types}에는 네 번의 실행이 나와요. 네 실행은 모두 2026년 10월 9일에 계획의 주소 집합에 대해 했어요. 그 주소 집합에서는 모든 토큰의 허락을 셌어요. 처음 후보 $20{,}000$개는 공개 접속 주소(블록체인 노드에 질문을 보내는 주소) \url{https://arb1.arbitrum.io/rpc}에서 읽었어요. 그러자 이 접속 주소가 요청 속도를 제한했어요. 그래서 나머지 $9{,}022$개는 같은 블록에서 \url{https://arbitrum-one.public.blastapi.io}로 읽었어요. 뒤의 실행들도 이 접속 주소를 썼어요.

\begin{table}[!htb]
\centering
\caption[\texttt{eth\_getCode}로 읽은 계정 종류.]{아비트럼 원에서 \texttt{eth\_getCode}로 읽은 계정 종류예요. 실행마다 한 블록에서 읽었어요. 계약은 코드가 비어 있지 않고, 그 코드가 위임 표시도 아닌 주소예요. 코드 없음과 EIP-7702(위임 표시)는 둘 다 외부 소유 계정(개인 키로 움직이는 계정)이에요. 주소 집합은 계획의 것이에요(어떤 토큰의 허락이든 셈해요). 수정 코호트는 이 집합들의 부분 집합이에요. W0 실행은 코호트 후보가 아닌 스펜더를 읽었어요. W1 실행은 앞의 어느 실행에서도 종류를 나누지 않은 주소를 읽었어요. 출처는 \texttt{data/1-raw-blockchain-extracts/cohort/}와 \texttt{gmx/}에서 결과 파일 옆에 있는 JSON 파일이에요.}
\label{tab:account-types}
\fitwidth{%
\begin{tabular}{lrrrrrr}
\toprule
주소 집합 & 블록 & 읽은 시각(UTC) & 주소 수 & 계약 & 코드 없음 & EIP-7702 \\
\midrule
코호트 후보: 2023년 10월 -- 2026년 6월에 주인 $\geq 3$명에게서 0이 아닌 허락을 받음 & 513{,}216{,}506 & 2026년 10월 9일, 13:55 & 29{,}022 & 27{,}844 & 1{,}121 & 57 \\
2023년 10월 -- 2026년 6월에 닫기를 $\geq 3$번 한 GMX~V2 계정 & 513{,}219{,}999 & 2026년 10월 9일, 14:10 & 40{,}307 & 7{,}633 & 28{,}075 & 4{,}599 \\
$t_1$에 0보다 큰 허락을 가진 W0 스펜더(코호트 후보가 아닌 것) & 513{,}233{,}527 & 2026년 10월 9일, 15:08 & 8{,}224 & 7{,}792 & 413 & 19 \\
W1 스펜더와 W1에서 닫기를 $\geq 3$번 한 GMX~V2 계정(앞에서 읽지 않은 것) & 513{,}257{,}476 & 2026년 10월 9일, 16:52 & 1{,}153 & 350 & 568 & 235 \\
\bottomrule
\end{tabular}
}
\end{table}

매칭 코호트(모든 점수가 똑같이 점수를 매기는 계약 묶음)는 계획의 코호트의 일부이고, 첫 번째 실행에서 종류를 나누었어요. W0 스펜더는 첫 번째와 세 번째 실행에서 종류를 나누었어요. $t_1$에 다섯 토큰 가운데 하나의 0보다 큰 마지막 허락(가장 최근 허락)을 가진 스펜더는, 계획의 에고 자료에서도 그 허락을 가지고 있었기 때문이에요. W1 스펜더와 트레이더 코호트는 네 실행 모두에서 계정 종류를 가져와요. 그래서 수정 분석에서는 새로 읽지 않았어요. 모든 읽기는 2026년 10월 9일의 체인을 기준으로 해요. 그리고 주소는 그날 가지고 있던 코드로 종류가 나뉘어요.

\dissertationSubheading[sec:benchmark-run-logs]{벤치마크 실행 기록}

표~\ref{tab:benchmark-runs}에는 벤치마크(속도 재기)에서 시간을 잰 풀이 프로그램 호출마다 실제 걸린 시간(초)을 실행한 순서대로 적었어요. 그리고 그 평균과 표준편차(값이 퍼진 정도)도 함께 적었어요. 설정마다, 시간 잰 실행 다섯 번 전에 시간을 재지 않는 준비 실행을 한 번 했어요. 그리고 그래프와 설정마다 시간은 한 번만 쟀어요. 표~\ref{tab:robustness-damping}의 $d=0.85$ 행과 표~\ref{tab:benchmark-scaling}의 전체 코호트 행은 이 표의 첫 번째 덩어리를 그대로 다시 보여 줘요.

% Generated by experiments/scripts/export_latex.py -- do not edit by hand.
% Source: experiments/data/2-processed-tables-and-evaluations/revised-usd-12m/benchmark/benchmark.json
\begin{table}[htbp]
\centering
\caption[벤치마크에서 시간을 잰 풀이 프로그램 호출 전부.]{벤치마크에서 시간을 잰 풀이 프로그램(solver) 호출마다 실제 걸린 시간을 초 단위로 보여 줘요(매칭 코호트, 계약 $n=6{,}483$개). 실행한 순서대로 적었고, 평균과 표준편차도 함께 적었어요. 설정마다 시간을 재지 않는 준비 실행을 한 번 한 뒤에 시간 잰 실행을 했어요. 그래프와 설정마다 시간은 한 번씩만 쟀어요. 표~\ref{tab:robustness-damping}의 $d=0.85$ 행과 표~\ref{tab:benchmark-scaling}의 코호트 전체 행은 첫 번째 묶음을 다시 쓴 것이에요. 단계는 표~\ref{tab:benchmark-scaling}의 부분 표본이에요.}
\label{tab:benchmark-runs}
\fitwidth{%
\begin{tabular}{lrrrrrrr}
\toprule
방법 & 실행 1 & 실행 2 & 실행 3 & 실행 4 & 실행 5 & 평균 & 표준편차 \\
\midrule
\multicolumn{8}{l}{\emph{전체 그래프, $d=0.85$}} \\
EndorseRank & 2.799 & 2.591 & 2.562 & 2.751 & 2.582 & 2.657 & 0.109 \\
AWP & 7.870 & 7.925 & 8.795 & 9.308 & 8.750 & 8.530 & 0.618 \\
C-PR ($\lambda=0.5$) & 11.737 & 11.304 & 11.429 & 11.527 & 11.821 & 11.563 & 0.214 \\
S-PR & 11.145 & 10.580 & 11.319 & 10.743 & 10.628 & 10.883 & 0.330 \\
\midrule
\multicolumn{8}{l}{\emph{전체 그래프, $d=0.75$}} \\
EndorseRank & 2.179 & 1.881 & 1.859 & 1.959 & 2.054 & 1.986 & 0.132 \\
AWP & 6.566 & 5.925 & 6.196 & 5.832 & 5.260 & 5.956 & 0.483 \\
\midrule
\multicolumn{8}{l}{\emph{전체 그래프, $d=0.95$}} \\
EndorseRank & 7.560 & 7.353 & 8.068 & 8.929 & 8.654 & 8.113 & 0.679 \\
AWP & 21.964 & 21.172 & 20.230 & 20.452 & 20.206 & 20.805 & 0.757 \\
\midrule
\multicolumn{8}{l}{\emph{단계 $n=1{,}000$}} \\
EndorseRank & 0.530 & 0.563 & 0.496 & 0.466 & 0.435 & 0.498 & 0.050 \\
AWP & 1.305 & 1.305 & 1.298 & 1.294 & 1.334 & 1.307 & 0.016 \\
\midrule
\multicolumn{8}{l}{\emph{단계 $n=2{,}000$}} \\
EndorseRank & 0.308 & 0.321 & 0.330 & 0.319 & 0.314 & 0.318 & 0.008 \\
AWP & 2.493 & 2.353 & 2.349 & 2.358 & 2.361 & 2.383 & 0.062 \\
\midrule
\multicolumn{8}{l}{\emph{단계 $n=5{,}000$}} \\
EndorseRank & 1.740 & 1.619 & 1.596 & 1.715 & 1.689 & 1.672 & 0.062 \\
AWP & 6.711 & 6.760 & 6.461 & 6.558 & 6.610 & 6.620 & 0.119 \\
\bottomrule
\end{tabular}
}
\end{table}


\dissertationSubheading[sec:parameter-summary]{매개변수 요약}

표~\ref{tab:parameter-summary}에는 수정 분석의 기본 매개변수를 모아 두었어요. 모든 값은 버전을 매긴 설정 파일에서 읽었어요. 다만 토큰, 가격, 금액 기울기들, 상한은 수정과 함께 커밋한 입력 파일에서 읽었어요(\ref{sec:amount-revision}절).

\begin{table}[!htb]
\centering
\caption{평가에 쓴 기본 매개변수.}
\label{tab:parameter-summary}
\fitwidth{%
\begin{tabular}{lll}
\toprule
매개변수 & 값 & 하는 일 \\
\midrule
페이지랭크 감쇠 계수 $d$ & $0.85$; 바꿔 본 값 $0.75$, $0.95$ & 순간 이동 섞기 \\
수렴 허용 오차 $\varepsilon$ & $L_1$ 노름으로 $10^{-8}$ & 거듭제곱 반복 멈춤 \\
최대 반복 수 $I_{\max}$ & $300$ & 거듭제곱 반복 상한 \\
로지스틱 $k$ & $0.01$ & 로지스틱 감쇠 기울기(AWP, EndorseRank, C-PR의 송금 층) \\
로지스틱 $t_0$ & $180$일 & 로지스틱 감쇠 가운데 점(AWP, EndorseRank, C-PR의 송금 층) \\
값 변환 $b$ & 달러당 $\ln 3/m$, $m_T=170$, $m_A=100$ & AWP($m_T$)와 EndorseRank($m_A$)의 금액 무게 $V(z)$ \\
허락 상한 $U$ & $19{,}000$달러 & C-PR 허락 층; 받은 허락 금액 \\
토큰과 가격 & WETH, USDC, USDT0, ARB, WBTC; DefiLlama 하루 종가 & 이벤트와 달러 금액(표~\ref{tab:usd-tokens}) \\
재시작 분포 & $\propto X_u$; 균등 & AWP; EndorseRank, C-PR \\
결합 무게 $\lambda$ & 주된 값 $0.5$; $0.25$, $0.75$ & C-PR 허락 층의 무게, 계획에서 정함 \\
S-PR 재시작 벡터 & 송금 층 점 위의 C-PR ($\lambda=1$), 합이 1이 되게 다시 맞춤 & 보증 씨앗 순간 이동 \\
점수 기간 & 2025-07-01 -- 2026-06-30; 2025-04-01 -- 2026-03-31 & 같은 기간과 W1; W0. UTC, 양 끝 포함; 동결일이 로지스틱 감쇠의 기준점 \\
추출(스캔~A) & 2023-10-01 -- 2026-06-30 & 등록 분석은 이 기간 전체로 점수를 매김 \\
W0 동결 / 라벨 & 2026-03-31 / 2026-04-01 -- 2026-06-30 & 스펜더 코호트와 트레이더 코호트; 규칙~A \\
W1 동결 / 라벨 & 2026-06-30 / 2026-07-01 -- 2026-09-30 & 스펜더 코호트와 트레이더 코호트; 규칙~B, 허용 폭 $\delta=0.01$ \\
부트스트랩 다시 뽑기 / 시드 & $400$ / $42$; $2{,}000$ / $42$ & 짝지은 계약 부트스트랩; 트레이더 비교 \\
시간 잰 실행 / 준비 실행 & $5$ / $1$ & 실행 시간 재는 방법 \\
규모 단계 시드 & \texttt{contract-benchmark-v1} & SHA256으로 늘 똑같이 뽑는 부분 표본 \\
빅쿼리 예산 & USD~50 & 모든 질의 전에 확인 \\
\bottomrule
\end{tabular}
}
\end{table}

% !TEX root = ../main.tex
\section{SQL 틀}
\label{ch:appendix-sql}

이 부록은 두 분석의 그래프가 바탕으로 삼는 질의 틀 세 가지를 줄여서 보여 줘요. 전체 파일은 실험 폴더의 \texttt{sql/}에 있어요(부록~\ref{sec:pipeline-stages}). \texttt{\_\_DEST\_\_}처럼 밑줄 두 개로 감싼 자리 표시는 질의를 돌리기 전에 \texttt{scripts/run\_bq.py}가 채워 넣어요. 그리고 \texttt{@}로 시작하는 이름은 질의 매개변수(질의에 밖에서 넣어 주는 값)예요.

\dissertationSubheading[sec:sql-materialize]{로그를 미리 표로 만들어 두는 틀}

스캔~A와 스캔~B는 같은 질의(\texttt{sql/01\_materialize\_logs.sql})를 서로 다른 기간에 대해 돌려요. 이 질의는 공개 로그 표를 한 번만 읽어요. 그리고 로그마다 종류 꼬리표를 붙이고, 이 연구에 필요한 네 종류를 달마다 나눈 프로젝트 표에 남겨요. 토픽(topic, 로그에 붙는 꼬리표) 해시와 GMX 갈래, Aave 갈래를 줄여서 보이면 다음과 같아요:

\begin{verbatim}[baselinestretch=1]
CREATE TABLE `__DEST__`
PARTITION BY TIMESTAMP_TRUNC(block_timestamp, MONTH) CLUSTER BY kind AS
WITH tagged AS (
  SELECT *, CASE
      WHEN ARRAY_LENGTH(topics) = 3 AND topics[OFFSET(0)] = '0x8c5b...b925'
        THEN 'A'                                  -- ERC-20 Approval
      WHEN ARRAY_LENGTH(topics) = 3 AND topics[OFFSET(0)] = '0xddf2...b3ef'
        THEN 'T'                                  -- ERC-20 Transfer
      WHEN ... THEN 'G'                           -- GMX V2 PositionDecrease
      WHEN ... THEN 'V'                           -- Aave V3 pool events
    END AS kind
  FROM `bigquery-public-data.goog_blockchain_arbitrum_one_us.logs`
  WHERE block_timestamp >= TIMESTAMP(@start_ts)
    AND block_timestamp < TIMESTAMP(@end_ts))
SELECT block_timestamp, block_number, log_index, kind,
  FROM_HEX(SUBSTR(address, 3)) AS emitter,
  IF(kind IN ('A', 'T'), FROM_HEX(SUBSTR(topics[OFFSET(1)], 27)), NULL) AS a,
  IF(kind IN ('A', 'T'), FROM_HEX(SUBSTR(topics[OFFSET(2)], 27)), NULL) AS b,
  IF(kind IN ('A', 'T'), SAFE.FROM_HEX(SUBSTR(data, 3)), NULL) AS val -- ...
FROM tagged WHERE kind IS NOT NULL
\end{verbatim}

허락에서는 \texttt{a}가 주인이고 \texttt{b}가 스펜더예요. 송금에서는 \texttt{a}가 보낸 쪽이고 \texttt{b}가 받는 쪽이에요. 두 경우 모두 \texttt{emitter}는 토큰 계약이고, \texttt{val}은 32바이트 금액이에요. 그리고 주소는 20바이트 값으로 저장해요.

\dissertationSubheading[sec:sql-ego]{에고 네트워크 뽑아내기 틀}

\texttt{sql/03\_ego\_events.sql} 질의는 한 기간에 대해 미리 만들어 둔 로그 표를 읽어요. 그리고 한쪽 끝에 계획의 코호트에 든 계약이 있는 허락과 송금만 남겨요. 계획의 코호트는 \texttt{contract\_rep.cohort} 표에 들어 있어요. 허락 부분을 줄여서 보이면 다음과 같아요:

\begin{verbatim}[baselinestretch=1]
CREATE OR REPLACE TABLE `dissertation-bq.contract_rep.approvals___SUFFIX__`
PARTITION BY TIMESTAMP_TRUNC(block_timestamp, MONTH) AS
SELECT l.block_timestamp, l.block_number, l.log_index,
       l.emitter AS token_address, l.a AS owner, l.b AS spender,
       u256(l.val) AS value
FROM `dissertation-bq.contract_rep.__SRC__` AS l
WHERE l.kind = 'A'
  AND l.block_timestamp >= TIMESTAMP(@start_ts)
  AND l.block_timestamp < TIMESTAMP(@end_ts)
  AND (l.a IN (SELECT address FROM `dissertation-bq.contract_rep.cohort`)
       OR l.b IN (SELECT address FROM `dissertation-bq.contract_rep.cohort`));
-- transfers___SUFFIX__: the same with kind = 'T', from_address = a,
-- to_address = b and zero amounts dropped: AND LTRIM(TO_HEX(l.val), '0') != ''
\end{verbatim}

파일 맨 위에 정의한 함수 \texttt{u256}은 32바이트 금액을 원래 단위(소수점 조정을 하지 않은 가장 작은 단위)의 배정밀도 실수(double)로 읽어요. 금액이 0인 허락은 남겨요. 금액이 0인 허락은 허락을 0으로 만드는 일이기 때문이에요. 반면 금액이 0인 송금은 버려요. 수정 분석(등록 뒤에 금액을 달러로 바꿔 다시 한 분석)은 이 질의를 다시 돌리지 않아요. 그 대신 \texttt{sql/12\_usd\_ego.sql}이 이 표들에서 다섯 토큰의 이벤트를 남겨요. 남기는 것은 스캔~A에서 다섯 토큰에 대해 주인 세 명 이상이 허락한 계약이 한쪽 끝에 있는 이벤트예요. 그리고 송금마다 그 블록이 있는 UTC 날짜의 토큰 가격으로 금액을 매겨요. 그다음 \texttt{sql/21\_usd12m\_cohorts.sql}이 각 동결일(점수를 고정하는 날)의 코호트를 그 동결일의 점수 기간(점수를 계산하는 데 쓰는 기간)에 있는 허락에서 뽑아요.

\dissertationSubheading[sec:sql-pairs]{마지막 허락과 쌍 표를 만드는 틀}

수정 분석에서는 \texttt{sql/14\_usd\_anchor\_pairs.sql} 질의가 한 기준 시점(앵커)까지의 에고 이벤트를 두 그래프의 화살표 목록으로 바꿔요. 등록 분석(미리 등록한 그대로 한 분석)에서는 \texttt{sql/04\_anchor\_pairs.sql} 질의가 같은 일을 원래 단위로, $b=1$을 쓰고 상한 없이 했어요. 줄여서 보이면 다음과 같아요:

\begin{verbatim}[baselinestretch=1]
CREATE OR REPLACE TABLE `dissertation-bq.contract_rep.u_allow_triples___TAG__` AS
SELECT r.token_address, r.owner, r.spender, r.amount * p.price_usd AS usd,
       TIMESTAMP_DIFF(TIMESTAMP(@anchor_ts), r.block_timestamp, MICROSECOND)
         / 8.64e10 AS age_days
FROM (
  SELECT *, ROW_NUMBER() OVER (
           PARTITION BY token_address, owner, spender
           ORDER BY block_number DESC, log_index DESC) AS rn
  FROM `dissertation-bq.contract_rep.u_approvals`
  WHERE block_timestamp >= TIMESTAMP(@start_ts)
    AND block_timestamp <= TIMESTAMP(@anchor_ts)
) AS r
JOIN `dissertation-bq.contract_rep.u_prices` AS p
  ON p.token = r.token_address AND p.day = DATE(TIMESTAMP(@anchor_ts))
WHERE r.rn = 1 AND r.value > 0;

CREATE OR REPLACE TABLE `dissertation-bq.contract_rep.u_allow_pairs___TAG__` AS
SELECT owner, spender,
       SUM(sig(age_days) * vt(usd, @b_allowance)) AS w_er,  -- EndorseRank
       SUM(LEAST(usd, @cap_allowance)) AS w_raw,    -- C-PR allowance layer
       MAX(sig(age_days)) AS max_sig
FROM `dissertation-bq.contract_rep.u_allow_triples___TAG__`
GROUP BY owner, spender;
-- u_transfer_pairs___TAG__: per (from_address, to_address) up to the anchor,
-- w_awp = SUM(sig * vt(usd, @b_transfer)), w_raw = SUM(usd * sig), MAX(sig)
\end{verbatim}

파일 맨 위에 정의한 함수 \texttt{sig}와 \texttt{vt}는 식~\eqref{eq:logistic-decay}의 로지스틱 감쇠(오래된 일일수록 가볍게 세는 무게)와 식~\eqref{eq:value-transform}의 값 변환(금액이 아무리 커도 1을 넘지 않게 바꾸는 식)이에요. 로지스틱 감쇠에서는 $k=0.01$이고 $t_0=180$일이에요. 값 변환은 $bz=40$부터 값을 1로 정해 두어요. 그 자리부터는 배정밀도로 계산해도 값이 1과 같기 때문이에요. 창 함수(window function, 묶음마다 행에 차례를 매기는 SQL 기능)는 토큰, 주인, 스펜더마다 블록 번호와 로그 번호(log index)로 마지막 허락을 골라요. 이 마지막 허락은 금액이 0보다 클 때만 남겨요. 그래서 허락 취소(허락을 거두어들임)가 있으면 그 토큰은 그 쌍에서 빠져요. 그리고 마지막 허락에는 기준 시점 날의 가격으로 금액을 매겨요. 1달러당 기울기 $b_A=\ln 3/100$과 $b_T=\ln 3/170$, 그리고 상한 $U=19{,}000$달러는 \texttt{constants.json}에서 매개변수로 넘겨받아요. 쌍 표의 각 행에는 모든 방법의 무게가 함께 들어 있어요. 그래서 한 기준 시점의 모든 점수는 같은 행들에서 나와요. 같은 질의는 기울기를 10배로 키운 경우와 10분의 1로 줄인 경우의 엔도스랭크(EndorseRank) 무게와 AWP 무게도 써 둬요. 표~\ref{tab:robustness-slope}의 금액 기울기 점검은 이것으로 계산해요.

% !TEX root = ../main.tex
\section{페이지랭크와 순위 상관의 기본 성질}
\label{ch:appendix-theory}

\ref{ch:literature}--\ref{ch:results}장은 여기에 정리한 기본 성질에 기대고 있어요. 이 내용은 널리 알려진 표준적인 것이에요. 페이지랭크의 성질은 \textcite{page1999pagerank}의 연구와 \textcite{langville2006pagerank}의 연구를 따라요. 순위 상관(두 줄 세우기가 얼마나 비슷한지 나타내는 값)은 \textcite{spearman1904proof}의 연구와 \textcite{kendall1938rank}의 연구를 따라요.

\dissertationSubheading[sec:pr-existence]{페이지랭크의 존재와 유일성}

$P$를 $|V|$개의 점 위에 있는 행 확률 행렬(각 행의 수를 모두 더하면 1인 행렬)이라고 해요. 필요하면 식~\eqref{eq:transition}처럼 매달린 점(나가는 화살표가 없는 점)을 고친 뒤의 행렬이에요. $d\in(0,1)$일 때 식~\eqref{eq:google-matrix}의 구글 행렬은 다음과 같아요.
\[
G_{\mathrm{PR}}=d P^{\top}+\frac{1-d}{|V|}\mathbf{1}\mathbf{1}^{\top},
\]
이 행렬은 모든 칸이 양수이고, 열 확률 행렬(각 열의 수를 모두 더하면 1인 행렬)이에요. 양수 행렬에 대한 페론--프로베니우스 정리(Perron--Frobenius theorem)에 따르면, $G_{\mathrm{PR}}$에는 크기(절댓값)가 1인 고윳값(행렬을 곱해도 방향이 바뀌지 않는 벡터가 늘어나는 배수)이 하나뿐이에요. 그리고 그 고윳값에 맞는 고유 벡터 $\mathbf{r}$은 모든 칸이 0보다 커요. 우리는 이 벡터를 $\sum_i r_i=1$이 되도록 크기를 맞출 수 있어요. 따라서 감쇠 계수(화살표를 따라갈 확률)와 매달린 점 규칙을 정해 두면, 방향이 있는 모든 그래프는 페이지랭크 점수를 단 하나만 가져요\parencite{langville2006pagerank}.

이 논문의 재시작 벡터 가운데 세 개는 고르지 않아요. AWP의 재시작 벡터는 활발함에 비례해요. 그래서 아무것도 보내지 않은 주소에서는 0이에요. 활동 재시작(활발한 주소에서 더 자주 다시 출발하기)을 쓰는 엔도스랭크(EndorseRank)의 재시작 벡터는 아무것도 허락하지 않은 주소에서 0이에요. S-PR의 재시작 벡터는 허락 걷기의 점수를 정규화(합이 1이 되게 맞춤)한 것이에요. 그래서 송금 층의 점 가운데 허락 점수가 없는 점에서는 0이에요. 따라서 세 구글 행렬 가운데 어느 것도, 모든 칸이 양수인 행렬이 아니에요. 그래도 고정점(계산을 한 번 더 해도 바뀌지 않는 값)은 여전히 하나뿐이에요. 두 확률 벡터의 차이를 구하면 다시 출발 항이 서로 지워져요. 그리고 이 연산자는 두 벡터 사이의 $\|\cdot\|_1$ 거리를 많아야 $d$배로 만들어요. 하지만 일반적으로 고정점의 모든 칸이 0보다 크지는 않아요. 다시 출발 몫이 없고, 다시 출발하는 점에서 출발한 어떤 길로도 닿지 않는 점은 점수가 0으로 남기 때문이에요. 답이 하나뿐이라는 것은, 매개변수를 정해 두면 각 방법이 주어진 화살표 집합에 점수 벡터를 꼭 하나만 준다는 뜻이에요. 그래서 두 점수 벡터 $\mathbf{s}_{\mathrm{ER}}$, $\mathbf{s}_{\mathrm{AWP}}$ 사이의 차이는 모두 두 방법의 화살표와 재시작 분포(어디서 다시 출발하는지의 비율)에서 생겨요.

\dissertationSubheading[sec:pr-continuity]{화살표 무게에 대한 연속성}

$d\in(0,1)$을 고정하면, 페이지랭크는 $P$의 칸 값에 따라 연속적으로(끊기지 않고 조금씩) 변해요. 같은 점들 위에 두 번째 행 확률 행렬 $\tilde P$가 있고, 그 페이지랭크 벡터를 $\tilde{\mathbf{r}}$로 써요. 그러면 다음이 성립해요.
\[
\|\mathbf{r}-\tilde{\mathbf{r}}\|_1\le\frac{d}{1-d}\,\max_i\bigl\|P_{i\cdot}-\tilde P_{i\cdot}\bigr\|_1,
\]
이 상한(넘지 않는 값)은 $|V|$에 따라 달라지지 않아요\parencite{langville2006pagerank}.

$P$에 대한 연속성은 화살표 무게에 대한 연속성과 같지 않아요. 행 $P_{i\cdot}=w_{i\cdot}/o_i$는 나가는 무게의 합 $o_i$가 0보다 큰 동안에만 무게에 따라 연속적으로 움직여요. 어떤 점이 첫 번째 나가는 화살표를 얻거나 마지막 나가는 화살표를 잃으면, 그 점의 행은 갑자기 뛰어요. 고르게 퍼진 매달린 점의 행과, 화살표가 가리키는 이웃들에 몰린 행 사이를 오가요. 관련된 무게가 아무리 작아도 그래요. \ref{sec:pagerank-worked-example}절의 계산 예에서 무게가 $10^{-12}$인 화살표 $D\to A$를 더하면, $\mathbf{r}$이 $\|\cdot\|_1$ 노름(차이의 절댓값을 모두 더한 크기)으로 $0.37$만큼 움직여요. 그리고 $r_A$는 $0.1375$에서 $0.3202$로 올라가요.

감쇠 계수를 여러 값으로 바꿔 보는 실험은 연속성으로 설명할 수 있어요. 이 실험은 $P$를 그대로 두고 $d$만 매끄럽게 바꿔요. 그리고 $d<1$이면 페이지랭크 벡터는 $d$에 따라 연속적으로 변해요. 금액 기울기 점검은 일부만 설명할 수 있어요. 모든 $z>0$에 대해 $V(z)>0$이에요. 그래서 기울기 $b$를 바꿔도, 정확한 계산에서는 화살표가 하나도 더해지거나 빠지지 않아요. 하지만 컴퓨터의 배정밀도 계산에서는 아주 작은 금액만 실은 화살표의 무게가 반올림되어 0이 될 수 있어요. 그리고 $b$를 10배 바꾸는 것은 너무 큰 변화라서, 위의 상한은 이런 변화에 대해 별로 말해 주지 못해요. 그래서 그 점검이 안정적인지는 자료에 달려 있어요(표~\ref{tab:robustness-slope}). 실제 자료에서는 $d\in\{0.75,0.85,0.95\}$에 걸쳐 엔도스랭크의 지표 가족(비슷한 지표 묶음) 평균이 많아야 $0.0031$만큼, AWP의 지표 가족 평균이 많아야 $0.021$만큼 움직였어요(표~\ref{tab:robustness-damping}).

\dissertationSubheading[sec:pr-sparsity]{희소성과 반복 한 바퀴의 비용}

0이 아닌 화살표가 $m=|E|$개, 점이 $n=|V|$개 있으면, 듬성듬성한 곱 $P^{\top}\mathbf{x}$에는 $O(m)$번의 산술 연산이 들어요. 반복 한 바퀴마다 길이가 $n$인 벡터들도 새로 고쳐요. 화살표 목록으로 듬성듬성한 연산자를 만드는 일은 $O(m)$이에요. 그래서 $I$번 반복하는 거듭제곱 반복(같은 계산을 되풀이하기)은 $O(m+(m+n)I)$예요. 두 점수는 같은 풀이 프로그램을 써요. 실제로 잰 호출에서 AWP의 반복 한 바퀴에는 엔도스랭크의 반복 한 바퀴보다 약 $1.95$배의 시간이 걸려요. 이것은 화살표 수의 비율인 $2.21$보다 작아요. 그리고 AWP는 반복을 더 많이 해야 해요(\ref{sec:benchmark-results}절). 그래서 $m$도 $I$도 혼자서는 시간 비율을 정하지 못해요.

\dissertationSubheading[sec:spearman-properties]{스피어먼의 \texorpdfstring{$\rho$}{rho}가 가진 성질}

스피어먼의 $\rho$는 $[-1,1]$ 안에 있어요. $\rho=1$이면 완벽한 단조 일치(한쪽이 커지면 다른 쪽도 늘 커짐)를 뜻해요. $\rho=-1$이면 완벽하게 거꾸로인 일치를 뜻해요. 평가 처리 과정은 동점을 제대로 다루는 표준 구현으로 $\rho$와 $\tau_b$를 계산해요. $n$이 크면 $\rho$와 $\tau$는 보통 부호가 같지만, 크기는 달라요. 엔도스랭크나 AWP를 대리 지표 하나와 짝지은 14개 쌍에서(표~\ref{tab:alignment-transfer}부터~\ref{tab:alignment-sybil-stability}까지), 모든 $|\tau|$는 $0.05$를 넘어요. 비율 $\rho/\tau$는 약 $1.2$와 $1.5$ 사이에 있어요. 다만 들어오는 상대 비율에 대한 AWP는 예외로, $\tau=0.055$이고 $\rho=0.163$이에요. 그리고 둘 가운데 늘 $|\rho|$가 더 커요. 계수의 부호가 서로 반대인 쌍은 하나도 없어요.

\dissertationSubheading[sec:kendall-properties]{켄달의 \texorpdfstring{$\tau$}{tau}가 가진 성질}

이 논문의 점수들에서는 동점이 드물어요. 매칭 코호트(모든 점수가 똑같이 점수를 매기는 계약 묶음)에서 어떤 점수도 계약 쌍의 $3.4\%$보다 많이 동점을 만들지 않아요. 엔도스랭크와 AWP의 동점 비율은 \ref{sec:rank-divergence}절에 있어요. 식~\eqref{eq:kendall}에서 $\tau_b$의 분모는 동점 때문에만 $n_0$과 달라요. 그래서 엔도스랭크와 AWP를 비교할 때 이 분모는 $n_0$과 2퍼센트 미만으로 달라요. 그리고 \ref{sec:rank-divergence}절의 쌍 개수로 계산한 $\tau_a$는 보고한 $\tau_b$와 2퍼센트 안쪽으로만 달라요.

대리 지표에 있는 동점도 $\tau_b$에 천장(얻을 수 있는 가장 큰 값)을 만들어요. 들어오는 화살표 수(입차수)나 받은 허락 수(허락 입차수) 같은 정수 대리 지표는 많은 쌍을 동점으로 만들어요. 그래서 그런 대리 지표에 대해 어떤 점수가 다다를 수 있는 가장 큰 $\tau_b$는 동점 쌍의 비율이 클수록 낮아져요. 페이지랭크 점수들은 동점 쌍이 적어요. 그래서 한 대리 지표에서 두 페이지랭크 점수를 비교하면, 이 천장은 둘 사이에서 거의 다르지 않아요. 따라서 그런 대비는 이 천장 때문에 비뚤어지지 않아요. 하지만 페이지랭크와 원시 차수(화살표를 그냥 센 수) 사이의 대비나, 두 대리 지표 사이의 대비는 비뚤어질 수 있어요. 둘의 동점이 서로 다른 만큼 그래요(\ref{sec:statistical-tests}절).

% !TEX root = ../main.tex
\section{자료 설명}
\label{ch:appendix-data}

이 부록은 \ref{ch:results}장의 바탕이 되는, 뽑아낸 자료를 설명해요. 다루는 것은 다음과 같아요. 에고 네트워크(코호트 계약에 닿는 화살표만 모은 그래프)를 뽑아낸 자료의 양과 두 그래프(점과 화살표로 이은 그림)의 크기, 다섯 토큰과 그 가격과 허락(allowance, 토큰을 얼마까지 대신 써도 된다는 허락) 금액, 그리고 GMX~V2 스마트 계약(블록체인 위에서 저절로 실행되는 프로그램) 계정의 분기별 활동이에요.

\dissertationSubheading[sec:data-events]{뽑아낸 이벤트와 그래프}

스캔~A의 에고 추출, 곧 2023년 10월 1일부터 2026년 6월 30일까지의 에고 추출은 계약 $27{,}844$개로 이루어진 계획의 코호트(묶음)를 위해 했어요. 이 추출에는 모든 ERC-20 토큰의 허락 로그(계약이 남기는 기록) $389{,}866{,}882$개와 송금(transfer) 로그 $1{,}538{,}404{,}067$개가 들어 있어요. 허락 로그는 금액이 0인 것도 포함해요. 송금 로그는 금액이 0이 아닌 것만이에요. 로그마다 한쪽 끝에 그 코호트의 계약이 있어요. 수정 분석(등록 뒤에 금액을 달러로 바꿔 다시 한 분석)은 동결일(점수를 고정하는 날)마다, 그 점수 기간(점수를 계산하는 데 쓰는 기간)에 있는 다섯 토큰의 허락과 송금 가운데 한쪽 끝에 그 동결일의 코호트 계약이 있는 것을 남겨요(그 수는 \ref{sec:dataset}절에 있어요). W1의 7월부터 9월까지의 로그 가운데서는 허락 $12{,}200{,}682$개와 송금 $58{,}625{,}425$개를 남겨요. 남긴 이벤트로 만든 쌍 표(주소 쌍마다 한 행씩 모은 표)에는 $T_{\mathrm{obs}}$에서 허락 쌍 $3{,}810{,}253$개와 송금 쌍 $8{,}432{,}635$개가 있어요. $t_1$에서는 $4{,}117{,}103$개와 $9{,}198{,}965$개가 있어요. $T_{\mathrm{obs}}$에서 금액 무게가 배정밀도(double) 계산에서 0이 되는 쌍 $1{,}304$개와 $6{,}511$개, 그리고 그런 쌍만 닿는 주소(\ref{sec:dataset}절)는 엔도스랭크(EndorseRank)와 AWP(적응형 가중 페이지랭크)의 풀이 프로그램이 쓰는 그래프에 들어가지 않아요(표~\ref{tab:benchmark-runtime}).

\dissertationSubheading[sec:data-tokens]{토큰, 가격, 허락 금액}

다섯 토큰을 고른 규칙은 \ref{sec:amount-revision}절에 있어요. $T_{\mathrm{obs}}$의 점수 기간에서, 이 다섯 토큰은 계획의 코호트의 에고 표에 있는 모든 토큰의 행 가운데 $79.40\%$를 차지해요. 여기서 행은 그 기간의 송금, 그리고 그 기간 안에 정해졌고 $T_{\mathrm{obs}}$에 0보다 큰 마지막 허락(가장 최근 허락)이에요. 표~\ref{tab:usd-tokens}에는 토큰마다 주소, 소수 자릿수, 가격 피드(블록체인에 올라오는 가격 정보), 에고 표에 있던 행과 남긴 행, 남긴 송금의 달러(USD) 금액 합, 그리고 하루 가격의 바탕이 된 자료가 있어요. 그 표의 캡션은 빠진 행을 세요. 빠진 행은 다른 토큰의 행, 그리고 다섯 토큰의 행 가운데 수정 코호트(수정 분석의 코호트)의 계약이 하나도 끼지 않은 행이에요.

% Generated by experiments/scripts/export_latex.py -- do not edit by hand.
% Source: experiments/data/0-usd-token-list-and-prices/listed_tokens.csv
% Source: experiments/data/2-processed-tables-and-evaluations/revised-usd-12m/describe/describe.json
% Source: experiments/data/0-usd-token-list-and-prices/chainlink_check.json
\begin{table}[htbp]
\centering
\caption[수정 분석의 다섯 토큰과 남긴 행, 그리고 가격.]{수정 분석에 쓴 다섯 토큰이에요. 토큰마다 다음을 적었어요. 먼저 체인링크(Chainlink) 가격 피드(시장 위험 등급)를 적었어요. 계획의 코호트가 가진 에고 표(에고 네트워크를 담은 표)에 $T_{\mathrm{obs}}$의 점수 기간(2025년 7월 1일부터 2026년 6월 30일까지) 동안 그 토큰의 행(허락과 송금)이 몇 개 있는지, 그리고 수정 분석이 그 가운데 몇 개를 남겼는지도 적었어요. 남긴 송금의 달러 거래 금액도 적었어요. 가격 시계열에서 두 번째 DefiLlama 키로 채운 날, 전날 값으로 채운 날, 하루짜리 오류로 보고 바꾼 날의 수도 적었어요(수정 분석의 기간들이 덮는 2025년 4월 1일부터 2026년 9월 30일까지의 548일 가운데). 마지막으로, 그 날들 가운데 점검 날짜 20개에 블록체인에서 읽은 Chainlink 피드와 DefiLlama 가격 사이 절대 차이의 중앙값과 가장 큰 값을 퍼센트로 적었어요. 계획의 코호트가 가진 에고 표에는 그 점수 기간에 허락 135{,}799{,}669건과 송금 580{,}802{,}372건이 있어요. 이 가운데 허락 49{,}138{,}325건과 송금 117{,}646{,}035건은 다른 토큰의 것이라서 뺐어요(그 점수 기간과 W1의 에고 표에 있는 다른 토큰은 44{,}119개예요). 다섯 토큰의 허락 674{,}782건과 송금 20{,}434{,}298건은 양 끝 어느 쪽도 수정 분석의 $T_{\mathrm{obs}}$ 코호트에 든 계약이 아니어서 뺐어요. 남긴 송금 가운데 그 토큰의 첫 가격보다 앞선 것은 0건이에요. $T_{\mathrm{obs}}$에서, 점수 기간 안에 정해진 0보다 큰 마지막 허락 4{,}185{,}287개 가운데 2{,}084{,}669개는 무제한 허락이에요. 그리고 2{,}331{,}420개는 C-PR 허락 층의 상한 $U=19{,}000$달러에 닿아요.}
\label{tab:usd-tokens}
\fitwidth{%
\begin{tabular}{llrlrrrrrrr}
\toprule
토큰 & 주소 & 소수 자릿수 & Chainlink 피드 & 행 수, 계획 & 행 수, 남김 & 송금 USD & 2번째 키 & 채움 & 오류 & CL 차이 (\%) \\
\midrule
WETH & \texttt{0x82af\ldots{}bab1} & 18 & ETH / USD (낮음) & 201{,}126{,}112 & 191{,}667{,}668 & $1.9\times10^{11}$ & 0 & 0 & 0 & 0.04 / 0.14 \\
USDC & \texttt{0xaf88\ldots{}5831} & 6 & USDC / USD (낮음) & 186{,}761{,}761 & 180{,}274{,}238 & $4.9\times10^{11}$ & 0 & 0 & 0 & 0.01 / 0.02 \\
USDT0 & \texttt{0xfd08\ldots{}cbb9} & 6 & USDT / USD (낮음) & 95{,}067{,}280 & 92{,}700{,}928 & $1.0\times10^{11}$ & 0 & 0 & 0 & 0.01 / 0.15 \\
WBTC & \texttt{0x2f2a\ldots{}5b0f} & 8 & WBTC / USD (낮음) & 41{,}730{,}499 & 40{,}610{,}808 & $7.7\times10^{10}$ & 0 & 0 & 0 & 0.05 / 0.24 \\
ARB & \texttt{0x912c\ldots{}6548} & 18 & ARB / USD (낮음) & 25{,}132{,}029 & 23{,}454{,}959 & $1.0\times10^{10}$ & 0 & 0 & 0 & 0.06 / 0.24 \\
\bottomrule
\end{tabular}
}
\end{table}


아비트럼 원(Arbitrum One) 주소로 찾은 DefiLlama의 가격 기록은 WETH와 WBTC는 2024년 8월에, USDT0는 2025년 1월에 시작해요. 이것은 첫 점수 기간이 열리기 전이에요. 그래서 점수 기간들이 덮는 $548$일은 모두 그 토큰 자신의 주소로 찾은 가격을 써요. 다섯 토큰 가운데 어느 것에도 전날 값으로 채운 날이나, 자료 오류로 보고 바꾼 날이 없어요. 그리고 남긴 송금 가운데 그 토큰의 첫 가격보다 앞선 것은 없어요. 토큰마다 이 날들 가운데 $20$일에 체인 위에서 읽은 Chainlink 가격 피드와 비교하면, DefiLlama 가격은 토큰마다 중앙값으로 $0.010\%$에서 $0.056\%$만큼 달라요. 가장 크게 다른 것은 WBTC로, $0.24\%$예요.

무제한 허락(금액 한도가 없는 허락)은 흔해요. $T_{\mathrm{obs}}$ 시점의 0보다 큰 마지막 허락 $4{,}185{,}287$개 가운데 $2{,}084{,}669$개($49.8\%$)는 원래 단위(소수점 조정을 하지 않은 가장 작은 단위)로 $1.15\times10^{77}$ 이상이에요. 이것은 $2^{256}-1$에 가까워요. 그리고 무제한 허락을 포함해 $2{,}331{,}420$개($55.7\%$)는 값이 상한 $U$(이보다 큰 금액은 이 값으로 자름) 이상이에요. 이 허락 $4{,}185{,}287$개는 주인(토큰 주인) $2{,}508{,}872$명의 것이에요. 그 가운데 $1{,}389{,}569$명은 무제한 허락을 적어도 하나 해 두었어요. 그리고 $256{,}803$명은 무제한 허락과 유한한 허락을 둘 다 해 두었어요. 등록 분석(미리 등록한 그대로 한 분석)은 모든 토큰을 원래 단위 그대로 두었어요. 그 분석에서는 에고 송금의 $99.52\%$와, $T_{\mathrm{obs}}$ 시점의 마지막 허락 $28{,}069{,}418$개 가운데 $99.78\%$가 원래 단위로 10 이상이었어요. 그 범위에서는 $b=1$일 때 $V$가 1과 $10^{-4}$ 안쪽으로만 차이 나요. 그리고 그 분석에서는 0보다 큰 마지막 허락을 해 둔 주인 $13{,}029{,}306$명 가운데 $1{,}532{,}301$명이 무제한 허락과 그보다 작은 허락을 둘 다 해 두었어요.

\dissertationSubheading[sec:data-gmx]{분기별 GMX~V2 계약 트레이더}

트레이더(거래하는 사람) 코호트는 계약 코드를 가진 GMX~V2 계정들이에요(\ref{sub:cohorts}절). 표~\ref{tab:gmx-quarters}는 이 계정들이 포지션을 닫은 횟수를 달력 분기마다 세요. 포지션 닫기는 이 서비스(프로토콜)의 이벤트 발신기(이벤트를 내보내는 계약)가 남기는 \texttt{PositionDecrease} 이벤트예요\parencite{gmxnddocs}. 표는 그 가운데 청산(손실이 너무 커져 거래가 강제로 끝나는 일)의 수도 함께 세요. 또 그 분기에 포지션을 적어도 하나 닫은 계약 계정 수와, 적어도 세 개 닫은 계약 계정 수도 세요.

\begin{table}[!htb]
\centering
\caption[계약 계정의 분기별 GMX~V2 포지션 닫기.]{계약 계정의 분기별 GMX~V2 포지션 닫기예요. 2023년 10월부터 2026년 6월까지: 이 기간에 포지션을 적어도 세 번 닫은 GMX~V2 계정 $40{,}307$개 가운데 계약 $7{,}633$개가 닫은 것이에요(표~\ref{tab:account-types}). 2026년 7월부터 9월까지(W1의 라벨 기간): 스캔~A의 계약 계정과 W1 실행의 계약 계정이 닫은 것이에요. 청산: 주문 종류가 청산이거나, 청산 처리기를 거쳐 나온 닫기예요. 출처: \texttt{data/2-processed-tables-and-evaluations/gmx/}에 있는 해독한 닫기 기록이에요.}
\label{tab:gmx-quarters}
\fitwidth{%
\begin{tabular}{lrrrr}
\toprule
분기 & 닫기 & 청산 & 닫기가 있는 계정 & 닫기가 $\geq 3$번인 계정 \\
\midrule
2023년 4분기 & 1{,}119 & 104 & 180 & 105 \\
2024년 1분기 & 7{,}956 & 692 & 1{,}247 & 961 \\
2024년 2분기 & 11{,}509 & 589 & 2{,}120 & 1{,}165 \\
2024년 3분기 & 7{,}947 & 456 & 1{,}807 & 724 \\
2024년 4분기 & 7{,}343 & 904 & 1{,}321 & 698 \\
2025년 1분기 & 6{,}929 & 498 & 497 & 273 \\
2025년 2분기 & 12{,}080 & 379 & 368 & 250 \\
2025년 3분기 & 28{,}753 & 897 & 1{,}750 & 1{,}567 \\
2025년 4분기 & 40{,}240 & 1{,}596 & 1{,}915 & 1{,}744 \\
2026년 1분기 & 6{,}026 & 233 & 269 & 171 \\
2026년 2분기 & 3{,}783 & 125 & 169 & 127 \\
\midrule
2026년 3분기 (W1) & 6{,}805 & 182 & 165 & 131 \\
\midrule
2023년 10월 -- 2026년 6월 & 133{,}685 & 6{,}473 & & \\
\bottomrule
\end{tabular}
}
\end{table}

계약 계정은 2025년 하반기에 포지션을 가장 많이 닫았고, 2026년에는 훨씬 적게 닫았어요. 트레이더 코호트에 들려면, 계정이 동결일에 그래프의 점이기도 해야 해요. 2026년 4월부터 6월까지 적어도 세 번 닫은 계약 계정 $127$개 가운데 $87$개가 W0의 트레이더 코호트를 이뤄요. 그리고 7월부터 9월까지의 $131$개 가운데 $70$개가 W1의 트레이더 코호트를 이뤄요.

% !TEX root = ../main.tex
\section{엔도스랭크와 AWP를 나란히 놓고 보기}
\label{ch:appendix-er-awp}

\ref{ch:results}장은 이 논문의 모든 점수를 함께 보고해요. 그리고 각 페이지랭크(PageRank)를, 그 페이지랭크가 매끄럽게 다듬는(평활하는) 원시 차수(화살표를 그냥 센 수)와 견주어 읽어요. 이 부록은 한 층만 걷는 두 점수인 엔도스랭크(EndorseRank)와 AWP(적응형 가중 페이지랭크)를 바로 견준 결과를 모아요. 같은 기간에서의 지표 가족(비슷한 지표 묶음) 평균과 짝지은 차이, 대리 지표(대신 재는 값) 하나하나에 대한 계수, 그리고 기간 밖에서의 짝지은 차이예요. 표~\ref{tab:endorserank-awp-holdout}의 트레이더(거래하는 사람) 행과 끝에 있는 격자 표를 빼면, 이 비교는 어느 것도 미리 정한 것이 아니고, 어느 것도 이 논문의 주장(C)을 담고 있지 않아요. 그 트레이더 행과 격자 표의 쌍은 W1과 함께 등록했고, 여기서는 재분석으로 고친 점수에서 계산해요. 두 점수는 서로 다른 화살표(엣지) 위를 걷고, 서로 다른 구성개념(점수가 재려고 하는 생각)을 위해 만들었어요. 그래서 점검마다, 앞서는 점수는 거의 언제나 그 점검이 세는 화살표를 걷는 점수예요. 두 화살표 종류 어느 쪽으로도 만들지 않은 결과는 예외예요. 어느 점수도 그 결과를 위해 만들지 않았기 때문이에요. 그 결과는 \ref{sub:neutral-label-results}절에서 보고해요. 두 점수가 점수를 매기는 범위도 달라요. 매칭 코호트(모든 점수가 똑같이 점수를 매기는 계약 묶음)의 계약 가운데 $191$개에는 허락 화살표가 하나도 닿지 않고, $1{,}192$개에는 송금 화살표가 하나도 닿지 않아요. 기간 밖에서는 어느 점수도 자기가 다듬는 원시 차수를 이기지 못해요. 나중의 새 송금자(새로 송금을 보내온 주소)를 줄 세우는 데서 AWP는 송금 입차수(송금을 보내온 서로 다른 주소 수)보다 못해요. 그리고 나중의 새 허락 쌍을 줄 세우는 데서 엔도스랭크는 받은 허락 수(허락 입차수)보다 못해요. 고르게 다시 출발할 때도, 활동 재시작(활발한 주소에서 더 자주 다시 출발하기)을 쓸 때도 그래요(\ref{sec:holdout-results}절과 \ref{sub:endorserank-restarts}절; 라벨을 본 뒤에 계산했어요). 그러므로 엔도스랭크가 AWP보다 앞선다고 해서, 허락을 따라 퍼뜨리는 것(전파)이 허락한 주인 수에 무엇인가를 더한다는 것이 드러나지는 않아요.

\dissertationSubheading[sec:er-awp-same-window]{같은 기간}

표~\ref{tab:alignment-hybrid}은 두 점수의 지표 가족 평균을 구간(믿을 만한 범위)과 함께 보여 줘요. 표~\ref{tab:tie-sensitivity}의 마지막 행은 두 점수 사이의 일치를 보여 줘요. 두 모양은 서로 엇갈려요. AWP는 송금 가족을 가장 가깝게 따라가고, 허락 가족을 가장 덜 따라가요. 반면에 엔도스랭크는 허락 가족을 가장 가깝게 따라가요. 표~\ref{tab:tau-diff}은 계약을 다시 뽑은 같은 표본에서 짝지은 차이 다섯 개를 검사해요(\ref{sec:statistical-tests}절). 이 차이들에는 표의 행 순서대로 i부터 v까지 번호를 붙였어요.

% Generated by experiments/scripts/export_latex.py -- do not edit by hand.
% Source: experiments/data/2-processed-tables-and-evaluations/revised-usd-12m/same-window/eval_summary.json
\begin{table}[htbp]
\centering
\caption[같은 기간에서 엔도스랭크와 AWP 사이의 짝지은 대비.]{같은 기간에서 엔도스랭크와 AWP의 켄달의 $\tau$ 계수 사이 차이를 짝지은 부트스트랩으로 검사한 결과예요(매칭 코호트, 계약 $n=6{,}483$개). 지표 가족 항목에는 지표 가족 평균 $\tau$를 써요. $a$와 $b$에 같은 계약 다시 뽑기를 써요. 그래서 구간은 짝지은 차이에 대한 구간이에요. 구간이 0을 벗어나면, 5\% 수준에서 표본을 뽑을 때 생기는 흔들림만으로는 설명할 수 없는 차이라는 뜻이에요. 신뢰구간(CI)은 계약을 다시 뽑아 구한 95\% 백분위 부트스트랩 구간이에요(짝지은 다시 뽑기 400번, 시드 42).}
\label{tab:tau-diff}
\fitwidth{%
\begin{tabular}{lccccc}
\toprule
대비 ($a$ 빼기 $b$) & $\tau(a)$ & $\tau(b)$ & $\Delta\tau$ & $\Delta\tau$의 95\% 신뢰구간 & $P(\Delta\tau>0)$ \\
\midrule
EndorseRank 허락 빼기 EndorseRank 송금 & 0.560 & 0.196 & +0.363 & [0.341, 0.382] & 1.000 \\
AWP 송금 빼기 EndorseRank 송금 & 0.736 & 0.196 & +0.539 & [0.518, 0.560] & 1.000 \\
EndorseRank 허락 빼기 AWP 허락 & 0.560 & 0.226 & +0.333 & [0.311, 0.354] & 1.000 \\
AWP 시빌 안정성 빼기 EndorseRank 시빌 안정성 & 0.419 & 0.026 & +0.393 & [0.373, 0.412] & 1.000 \\
EndorseRank 받은 허락 수 빼기 EndorseRank 입차수 & 0.650 & 0.218 & +0.432 & [0.410, 0.452] & 1.000 \\
\bottomrule
\end{tabular}
}
\end{table}


다섯 구간은 모두 0보다 위에 있어요. AWP의 송금 일치도는 엔도스랭크의 송금 일치도보다 높아요(대비~ii, $+0.539$, $[0.518,0.560]$). 엔도스랭크의 허락 일치도는 AWP의 허락 일치도보다 높아요(대비~iii, $+0.333$, $[0.311,0.354]$). 그리고 AWP는 시빌 안정성 가족(시빌 공격, 곧 가짜 주소를 많이 만들어 속이는 공격을 생각해 고른 지표 묶음)에서 엔도스랭크보다 앞서요(대비~iv, $+0.393$, $[0.373,0.412]$). 나머지 두 대비는 엔도스랭크를 자기 자신과 견줘요. 엔도스랭크의 허락 일치도는 자기의 송금 일치도보다 높아요(대비~i, $+0.363$, $[0.341,0.382]$). 그리고 받은 허락 수와의 일치도는 들어오는 화살표 수(입차수)와의 일치도보다 높아요(대비~v, $+0.432$, $[0.410,0.452]$). 엔도스랭크는 계약 쌍의 $0.20\%$만 동점으로 둬요. 그래서 엔도스랭크 자신의 동점은 이 계수들 어느 것에도 천장(얻을 수 있는 가장 큰 값)을 씌우지 않아요. 그리고 대리 지표의 동점으로는 이만큼 큰 차이를 설명할 수 없어요(부록~\ref{sec:kendall-properties}). 허락 그래프 밖의 계약들에 0점을 주는 대신 외톨이 점(화살표가 하나도 없는 점)으로 남겨 두어도, 지표 가족 평균은 하나도 바뀌지 않아요(표~\ref{tab:tie-sensitivity}).

대비~ii와~iii의 방향이 서로 반대인 것은 구성개념에서 나와요. 각 점수는 자기 화살표로 만든 가족에서 앞서요. 그리고 가족마다 주된 대리 지표에는, 그에 맞는 페이지랭크가 다듬는 입차수가 들어 있어요. 그래서 이 일치는 일부가 만든 방식 때문에 저절로 생기고, 각 점수의 의도한 구성개념을 점검해요. 어느 차이도 어느 점수가 평판을 더 잘 재는지는 말해 주지 않아요.

\dissertationSubheading[sec:proxy-family-detail]{대리 지표 하나하나 보기}

지표 가족 평균 하나는 대리 지표마다 구한 계수 여러 개를 한데 줄인 값이에요. 표~\ref{tab:alignment-transfer}부터~\ref{tab:alignment-sybil-stability}까지는 엔도스랭크와 AWP의 계수를 대리 지표 하나씩 보고해요. 스피어먼의 로 $\rho$(두 줄 세우기가 얼마나 비슷한지 나타내는 수), 켄달의 타우 $\tau$(두 줄 세우기가 얼마나 같은 순서인지 나타내는 수), 그리고 $\tau$의 95\% 구간이에요. 송금 대리 지표는 AWP를 만든 사람들이 AWP를 평가할 때 쓴 바로 그 지표예요\parencite{do2023awp}.

\dissertationSubsubheading[sub:transfer-family]{송금 가족}

표~\ref{tab:alignment-transfer}은 두 점수와 입차수, 그리고 받은 양의 합 사이의 계수를 보여 줘요.

% Generated by experiments/scripts/export_latex.py -- do not edit by hand.
% Source: experiments/data/2-processed-tables-and-evaluations/revised-usd-12m/same-window/eval_summary.json
\begin{table}[htbp]
\centering
\caption[송금 대리 지표에 대한 엔도스랭크와 AWP의 분야 일치도.]{엔도스랭크와 AWP가 송금 대리 지표(각 계약의 입차수와 받은 양의 합)와 얼마나 잘 맞는지 나타낸 분야 일치도예요. 같은 기간에서 쟀어요(매칭 코호트, 계약 $n=6{,}483$개). 점수의 순위와 대리 지표의 순위 사이의 켄달의 $\tau$와 스피어먼의 $\rho$를 보여 줘요. 신뢰구간(CI)은 계약을 다시 뽑아 구한 95\% 백분위 부트스트랩 구간이에요(짝지은 다시 뽑기 400번, 시드 42).}
\label{tab:alignment-transfer}
\fitwidth{%
\begin{tabular}{lccc}
\toprule
대리 지표 & 스피어먼의 $\rho$ & 켄달의 $\tau$ & $\tau$의 95\% 신뢰구간 \\
\midrule
\multicolumn{4}{l}{\textit{EndorseRank}} \\
입차수 & 0.302 & 0.218 & [0.199, 0.238] \\
받은 양의 합 (들어온 양을 모두 더함) & 0.256 & 0.175 & [0.157, 0.194] \\
\addlinespace
\multicolumn{4}{l}{\textit{AWP}} \\
입차수 & 0.883 & 0.742 & [0.732, 0.752] \\
받은 양의 합 (들어온 양을 모두 더함) & 0.887 & 0.730 & [0.720, 0.740] \\
\bottomrule
\end{tabular}
}
\end{table}


AWP는 입차수와 맞는 것과 거의 같은 만큼 받은 양의 합과도 잘 맞아요(표~\ref{tab:alignment-transfer}). 그 까닭은 \ref{sec:same-window-alignment}절에 있어요. 엔도스랭크의 송금 계수의 구간은 0과도 겹치지 않고, AWP의 구간과도 겹치지 않아요. 이 계수들은 주로, 많은 주인이 허락한 계약이 코호트의 나머지 계약보다 더 많은 보낸 쪽(송신자)에게서 송금을 받는다는 것을 보여 줘요.

\dissertationSubsubheading[sub:allowance-family]{허락 가족}

표~\ref{tab:alignment-allowance}은 두 점수와 받은 허락 수, 그리고 받은 허락 금액 사이의 계수를 보여 줘요.

% Generated by experiments/scripts/export_latex.py -- do not edit by hand.
% Source: experiments/data/2-processed-tables-and-evaluations/revised-usd-12m/same-window/eval_summary.json
\begin{table}[htbp]
\centering
\caption[허락 대리 지표에 대한 엔도스랭크와 AWP의 분야 일치도.]{엔도스랭크와 AWP가 허락 대리 지표(각 스펜더 계약의 받은 허락 수와 받은 허락의 양)와 얼마나 잘 맞는지 나타낸 분야 일치도예요. 같은 기간에서 쟀어요(매칭 코호트, 계약 $n=6{,}483$개). 점수의 순위와 대리 지표의 순위 사이의 켄달의 $\tau$와 스피어먼의 $\rho$를 보여 줘요. 신뢰구간(CI)은 계약을 다시 뽑아 구한 95\% 백분위 부트스트랩 구간이에요(짝지은 다시 뽑기 400번, 시드 42).}
\label{tab:alignment-allowance}
\fitwidth{%
\begin{tabular}{lccc}
\toprule
대리 지표 & 스피어먼의 $\rho$ & 켄달의 $\tau$ & $\tau$의 95\% 신뢰구간 \\
\midrule
\multicolumn{4}{l}{\textit{EndorseRank}} \\
받은 허락 수 & 0.804 & 0.650 & [0.639, 0.661] \\
받은 허락의 양 (합) & 0.628 & 0.469 & [0.454, 0.484] \\
\addlinespace
\multicolumn{4}{l}{\textit{AWP}} \\
받은 허락 수 & 0.320 & 0.242 & [0.223, 0.262] \\
받은 허락의 양 (합) & 0.294 & 0.211 & [0.193, 0.229] \\
\bottomrule
\end{tabular}
}
\end{table}


엔도스랭크는 받은 허락 금액보다 받은 허락 수를 더 가깝게 따라가요. 그 까닭은 \ref{sec:same-window-alignment}절에 있어요. AWP의 계수는 0보다 위에 있어요. 코호트의 계약 가운데, 많은 주인이 허락한 계약은 송금 그래프에서도 중심에 있는 경향이 있어요.

\dissertationSubsubheading[sub:sybil-family]{시빌 보정 안정성 가족}

표~\ref{tab:alignment-sybil-stability}는 두 점수와 들어오는 상대 비율, 송금 활동 기간, 활동한 달 수 사이의 계수를 보여 줘요.

% Generated by experiments/scripts/export_latex.py -- do not edit by hand.
% Source: experiments/data/2-processed-tables-and-evaluations/revised-usd-12m/same-window/eval_summary.json
\begin{table}[htbp]
\centering
\caption[시빌 보정 안정성 대리 지표에 대한 엔도스랭크와 AWP의 분야 일치도.]{엔도스랭크와 AWP가 송금 활동에서 구한 시빌 보정 안정성 대리 지표와 얼마나 잘 맞는지 나타낸 분야 일치도예요. 같은 기간에서 쟀어요(매칭 코호트, 계약 $n=6{,}483$개). 점수의 순위와 대리 지표의 순위 사이의 켄달의 $\tau$와 스피어먼의 $\rho$를 보여 줘요. 신뢰구간(CI)은 계약을 다시 뽑아 구한 95\% 백분위 부트스트랩 구간이에요(짝지은 다시 뽑기 400번, 시드 42).}
\label{tab:alignment-sybil-stability}
\fitwidth{%
\begin{tabular}{lccc}
\toprule
대리 지표 & 스피어먼의 $\rho$ & 켄달의 $\tau$ & $\tau$의 95\% 신뢰구간 \\
\midrule
\multicolumn{4}{l}{\textit{EndorseRank}} \\
들어오는 상대 비율 & -0.290 & -0.197 & [-0.213, -0.180] \\
송금 활동 기간 (일) & 0.199 & 0.133 & [0.114, 0.153] \\
활동한 달 수 & 0.203 & 0.143 & [0.124, 0.164] \\
\addlinespace
\multicolumn{4}{l}{\textit{AWP}} \\
들어오는 상대 비율 & 0.163 & 0.055 & [0.033, 0.079] \\
송금 활동 기간 (일) & 0.763 & 0.588 & [0.575, 0.600] \\
활동한 달 수 & 0.768 & 0.616 & [0.603, 0.629] \\
\bottomrule
\end{tabular}
}
\end{table}


AWP는 송금 활동 기간과 활동한 달 수에서 엔도스랭크보다 앞서요. 기간 전체에 퍼진 활동에 점수를 더 주는 송금 점수라면 그럴 것이라고 예상할 수 있어요. 지표 가족 평균에서의 짝지은 차이는 0을 벗어나요(표~\ref{tab:tau-diff}, 대비~iv). 들어오는 상대 비율에서는 두 계수의 부호가 서로 반대예요. 엔도스랭크는 음수이고, AWP는 양수예요. 이 비율은 같은 보낸 쪽이 다시 오면 내려가요. 그래서 많은 주인이 허락한 계약은, 이 잣대로 보면 보낸 쪽이 다시 찾아오는 계약이에요. AWP의 계수가 양수인 것은 점수를 매기는 범위 때문이에요. 송금 그래프 밖에 있는 계약 $1{,}192$개는 비율도 0이고 AWP 점수도 0이에요. 그래서 이 계약들이 두 줄 세우기의 맨 아래를 함께 채워요. 이 가족은 각 점수가 시간에 걸친 송금 정보를 얼마나 담고 있는지를 보여 줘요.

\dissertationSubheading[sec:er-awp-holdout]{기간 밖}

표~\ref{tab:endorserank-awp-holdout}은 두 홀드아웃(나중 기간을 떼어 두고 맞혀 보는 시험) 기간에서 두 점수 사이의 짝지은 차이를 모아요. 모든 행에서 $a$는 엔도스랭크이고 $b$는 AWP예요. 스펜더(spender, 허락을 받아 대신 쓰는 쪽) 행은 라벨을 본 뒤에 계산했어요.

% Generated by experiments/scripts/export_latex.py -- do not edit by hand.
% Source: experiments/data/2-processed-tables-and-evaluations/revised-usd-12m/holdout-w0/eval_summary.json
% Source: experiments/data/2-processed-tables-and-evaluations/revised-usd-12m/supplementary/supplementary.json
% Source: experiments/data/2-processed-tables-and-evaluations/revised-usd-12m/window-w1/eval_summary.json
\begin{table}[htbp]
\centering
\caption[기간 밖에서 엔도스랭크와 AWP 사이의 짝지은 부트스트랩 차이.]{기간 밖에서 엔도스랭크와 AWP 사이의 차이를 짝지은 부트스트랩으로 구했어요(W0 스펜더 $n=7{,}914$, W0 트레이더 $n=87$, W1 스펜더 $n=7{,}960$, W1 트레이더 $n=70$). 이 가운데 어느 것도 규칙 A나 규칙 B에 들어가지 않아요. 스펜더 행은 라벨을 본 뒤에 계산했어요. 트레이더 행은 어느 화살표 종류로도 만들지 않은 결과에서 한 비교 가운데, 엔도스랭크에서 AWP를 뺀 쌍이에요(표~\ref{tab:neutral-label-contrasts}와~\ref{tab:neutral-label-grid}). W1 스펜더 행은 W1 평가가 끝난 뒤에, 저장해 둔 W1 점수로 계산했어요. 다시 뽑은 표본은 표~\ref{tab:fresh-holdout}와 같은 것을 썼어요. 각 묶음 안에서 $a$와 $b$는 같은 계약 다시 뽑은 표본을 함께 써요(스펜더는 400번, 트레이더는 2{,}000번, 시드 42).}
\label{tab:endorserank-awp-holdout}
\fitwidth{%
\begin{tabular}{llccccc}
\toprule
대비 ($a$ 빼기 $b$) & 라벨 & $\tau(a)$ & $\tau(b)$ & $\Delta\tau$ & $\Delta\tau$의 95\% 신뢰구간 & $P(\Delta\tau>0)$ \\
\midrule
\multicolumn{7}{l}{\emph{W0, 스펜더 ($n=7{,}914$)}} \\
EndorseRank 빼기 AWP & 새 허락 쌍 & 0.303 & 0.196 & +0.107 & [0.086, 0.130] & 1.000 \\
EndorseRank 빼기 AWP & 새 송금자 & 0.272 & 0.409 & -0.138 & [-0.151, -0.123] & 0.000 \\
\midrule
\multicolumn{7}{l}{\emph{W0, 세 번 이상 닫은 트레이더 ($n=87$)}} \\
EndorseRank 빼기 AWP & 청산 없이 닫은 비율 & 0.236 & 0.254 & -0.019 & [-0.173, 0.144] & 0.401 \\
\midrule
\multicolumn{7}{l}{\emph{W1, 스펜더 ($n=7{,}960$)}} \\
EndorseRank 빼기 AWP & 새 허락 쌍 & 0.299 & 0.157 & +0.142 & [0.119, 0.164] & 1.000 \\
EndorseRank 빼기 AWP & 새 송금자 & 0.259 & 0.398 & -0.139 & [-0.154, -0.125] & 0.000 \\
\midrule
\multicolumn{7}{l}{\emph{W1, 세 번 이상 닫은 트레이더 ($n=70$)}} \\
EndorseRank 빼기 AWP & 청산 없이 닫은 비율 & 0.297 & 0.263 & +0.034 & [-0.067, 0.134] & 0.750 \\
\bottomrule
\end{tabular}
}
\end{table}


W0의 스펜더에서 엔도스랭크는 새 허락 쌍을 AWP보다 더 잘 줄 세워요. 짝지은 차이는 $+0.107$(구간 $[0.086,0.130]$)이에요. 새 송금자에서는 순서가 뒤집혀요. 차이는 $-0.138$($[-0.151,-0.123]$)이에요. 모든 라벨은 점수에 쓴 모든 이벤트보다 뒤에 와요. 첫 번째 차이는 점수를 매기는 범위에 달려 있어요. AWP는 동결일 전에 송금 화살표가 없던 스펜더 $1{,}342$개에 0점을 줘요. 이것은 대략 여섯 개에 하나꼴이에요. 이 스펜더들은 허락 화살표로만 순서가 정해져요. 그리고 이 스펜더들은 나머지 스펜더보다 새 허락 쌍을 더 자주 얻었어요($26.7\%$ 대 $20.9\%$). 다른 주소에게서 송금을 받은 적이 있는 스펜더 $6{,}489$개에서는 AWP가 두 라벨 모두에서 앞서요(표~\ref{tab:holdout-receiving}). 여기에는 짝지은 검사를 하지 않았어요.

7월부터 9월까지에서도 AWP는 송금 화살표가 없는 스펜더에 다시 0점을 줘요. 그리고 같은 모양이 되풀이돼요. 엔도스랭크는 새 허락 쌍에서 앞서고($+0.142$, $[0.119,0.164]$), AWP는 새 송금자에서 앞서요($-0.139$, $[-0.154,-0.125]$).

W0 트레이더에서는(표~\ref{tab:holdout-traders}) 이익을 낸 닫기의 수와 실현 이익(실제로 얻은 이익)에서 AWP의 구간만 0보다 위에 있어요. 이 두 라벨은 닫기 횟수가 많을수록 커져요. 두 비율은 그렇지 않은데, 여기서는 두 점수가 모두 음수예요. 그리고 손실이 아닌 비율에서만 구간이 0보다 아래에 있어요. 어느 점수도 이 계약 트레이더(GMX~V2에서 거래하는 계약)들을 성공한 닫기의 비율대로 줄 세우지 못해요.

\dissertationSubheading[sec:er-awp-neutral-grid]{두 기간의 모든 트레이더 라벨}

표~\ref{tab:neutral-label-grid}은 두 기간에서 트레이더 라벨 여섯 개 모두에 대해 허락 쪽과 송금 쪽을 맞세워요. 이 표의 쌍은 이 연구가 7월부터 9월까지의 로그를 뽑기 전에 W1과 함께 등록했어요(\ref{sub:neutral-label-protocol}절). 여기서는 등록 분석의 결과를 안 뒤에 그 쌍을 고친 점수로 계산해요. 그래서 이 표는 재분석이에요(부록~\ref{sec:map-registered}). 청산 없이 닫은 비율에서의 P와 Q는 규칙 R1과 R2에 들어가요. 그리고 두 비율에서의 P는 규칙 R3에 들어가요(주장~C6, \ref{sub:neutral-label-results}절). W1에서는 두 청산 라벨 모두에서 개수(P)가 허락 쪽에 유리하고, 구간이 0보다 위에 있어요. 반면 세 페이지랭크 쌍은 양수이지만, 구간이 0을 품어요. W0에서는 두 청산 라벨 모두에서 P가 양수이지만, 구간이 0을 품어요. 그리고 세 페이지랭크 쌍은 0보다 아래에 있고, 그 구간도 0을 품어요. W1의 이익 라벨 네 개에서는 모든 쌍이 송금 쪽에 유리해요. Q, C-PR 층, 엔도스랭크 쌍의 구간은 네 라벨 모두에서 0보다 아래에 있어요. P의 구간은 두 비율에서 0보다 아래에 있어요. W0에서는 이익 라벨의 모든 구간이 0을 품어요. 다만 실현 이익에서의 C-PR 층의 구간은 예외예요.

% Generated by experiments/scripts/export_latex.py -- do not edit by hand.
% Source: experiments/data/2-processed-tables-and-evaluations/revised-usd-12m/window-w1/eval_summary.json
\begin{table}[htbp]
\centering
\caption[두 기간의 모든 트레이더 라벨에서 허락 쪽 점수 빼기 송금 쪽 점수.]{허락 쪽 점수에서 송금 쪽 점수를 뺀 $\Delta\tau$와 95\% 짝지은 부트스트랩 구간이에요. 기간 W1(GMX~V2 계약 트레이더 $n=70$)과 W0($n=87$)의 트레이더 라벨 여섯 개 모두에서 구했어요. P는 받은 허락 수 빼기 송금 입차수예요. Q는 활동 재시작을 쓰는 엔도스랭크 빼기 AWP예요(화살표 무게와 다시 출발 규칙이 같음). C-PR은 $\lambda=1$ 빼기 $\lambda=0$이에요(같은 달러 방식으로 만듦). ER은 엔도스랭크 빼기 AWP예요. 이익을 낸 닫기 수와 실현 이익은 닫은 횟수가 많을수록 커져요. 다른 네 라벨은 그렇지 않아요. 신뢰구간(CI)은 계약을 다시 뽑아 구한 95\% 백분위 부트스트랩 구간이에요(짝지은 다시 뽑기 2{,}000번, 시드 42).}
\label{tab:neutral-label-grid}
\fitwidth{%
\begin{tabular}{lcccc}
\toprule
라벨 & P (개수) & Q (페이지랭크) & C-PR 층 & ER 빼기 AWP \\
\midrule
\multicolumn{5}{l}{\emph{기간 W1, 2026년 7월--9월 (트레이더 $n=70$)}} \\
청산 없이 닫은 비율 & \begin{tabular}[c]{@{}c@{}}0.141\\[-1pt]{\footnotesize [0.037, 0.240]}\end{tabular} & \begin{tabular}[c]{@{}c@{}}0.023\\[-1pt]{\footnotesize [-0.062, 0.113]}\end{tabular} & \begin{tabular}[c]{@{}c@{}}0.019\\[-1pt]{\footnotesize [-0.093, 0.128]}\end{tabular} & \begin{tabular}[c]{@{}c@{}}0.034\\[-1pt]{\footnotesize [-0.067, 0.134]}\end{tabular} \\
청산 없음 & \begin{tabular}[c]{@{}c@{}}0.141\\[-1pt]{\footnotesize [0.044, 0.232]}\end{tabular} & \begin{tabular}[c]{@{}c@{}}0.027\\[-1pt]{\footnotesize [-0.066, 0.129]}\end{tabular} & \begin{tabular}[c]{@{}c@{}}0.021\\[-1pt]{\footnotesize [-0.087, 0.132]}\end{tabular} & \begin{tabular}[c]{@{}c@{}}0.038\\[-1pt]{\footnotesize [-0.069, 0.149]}\end{tabular} \\
이익을 낸 닫기 수 & \begin{tabular}[c]{@{}c@{}}-0.115\\[-1pt]{\footnotesize [-0.257, 0.025]}\end{tabular} & \begin{tabular}[c]{@{}c@{}}-0.215\\[-1pt]{\footnotesize [-0.399, -0.042]}\end{tabular} & \begin{tabular}[c]{@{}c@{}}-0.213\\[-1pt]{\footnotesize [-0.402, -0.028]}\end{tabular} & \begin{tabular}[c]{@{}c@{}}-0.200\\[-1pt]{\footnotesize [-0.390, -0.023]}\end{tabular} \\
실현 이익 & \begin{tabular}[c]{@{}c@{}}-0.109\\[-1pt]{\footnotesize [-0.247, 0.032]}\end{tabular} & \begin{tabular}[c]{@{}c@{}}-0.299\\[-1pt]{\footnotesize [-0.457, -0.146]}\end{tabular} & \begin{tabular}[c]{@{}c@{}}-0.294\\[-1pt]{\footnotesize [-0.476, -0.116]}\end{tabular} & \begin{tabular}[c]{@{}c@{}}-0.292\\[-1pt]{\footnotesize [-0.460, -0.134]}\end{tabular} \\
이익을 낸 비율 & \begin{tabular}[c]{@{}c@{}}-0.162\\[-1pt]{\footnotesize [-0.295, -0.025]}\end{tabular} & \begin{tabular}[c]{@{}c@{}}-0.262\\[-1pt]{\footnotesize [-0.436, -0.100]}\end{tabular} & \begin{tabular}[c]{@{}c@{}}-0.266\\[-1pt]{\footnotesize [-0.461, -0.070]}\end{tabular} & \begin{tabular}[c]{@{}c@{}}-0.246\\[-1pt]{\footnotesize [-0.427, -0.073]}\end{tabular} \\
손실이 아닌 비율 & \begin{tabular}[c]{@{}c@{}}-0.222\\[-1pt]{\footnotesize [-0.328, -0.106]}\end{tabular} & \begin{tabular}[c]{@{}c@{}}-0.324\\[-1pt]{\footnotesize [-0.471, -0.188]}\end{tabular} & \begin{tabular}[c]{@{}c@{}}-0.316\\[-1pt]{\footnotesize [-0.481, -0.152]}\end{tabular} & \begin{tabular}[c]{@{}c@{}}-0.315\\[-1pt]{\footnotesize [-0.468, -0.176]}\end{tabular} \\
\midrule
\multicolumn{5}{l}{\emph{기간 W0, 2026년 4월--6월 (트레이더 $n=87$)}} \\
청산 없이 닫은 비율 & \begin{tabular}[c]{@{}c@{}}0.095\\[-1pt]{\footnotesize [-0.027, 0.234]}\end{tabular} & \begin{tabular}[c]{@{}c@{}}-0.012\\[-1pt]{\footnotesize [-0.153, 0.149]}\end{tabular} & \begin{tabular}[c]{@{}c@{}}-0.057\\[-1pt]{\footnotesize [-0.209, 0.092]}\end{tabular} & \begin{tabular}[c]{@{}c@{}}-0.019\\[-1pt]{\footnotesize [-0.173, 0.144]}\end{tabular} \\
청산 없음 & \begin{tabular}[c]{@{}c@{}}0.109\\[-1pt]{\footnotesize [-0.022, 0.257]}\end{tabular} & \begin{tabular}[c]{@{}c@{}}-0.005\\[-1pt]{\footnotesize [-0.162, 0.165]}\end{tabular} & \begin{tabular}[c]{@{}c@{}}-0.056\\[-1pt]{\footnotesize [-0.211, 0.098]}\end{tabular} & \begin{tabular}[c]{@{}c@{}}-0.007\\[-1pt]{\footnotesize [-0.170, 0.165]}\end{tabular} \\
이익을 낸 닫기 수 & \begin{tabular}[c]{@{}c@{}}-0.002\\[-1pt]{\footnotesize [-0.141, 0.129]}\end{tabular} & \begin{tabular}[c]{@{}c@{}}-0.027\\[-1pt]{\footnotesize [-0.170, 0.113]}\end{tabular} & \begin{tabular}[c]{@{}c@{}}-0.089\\[-1pt]{\footnotesize [-0.244, 0.070]}\end{tabular} & \begin{tabular}[c]{@{}c@{}}-0.046\\[-1pt]{\footnotesize [-0.205, 0.113]}\end{tabular} \\
실현 이익 & \begin{tabular}[c]{@{}c@{}}-0.026\\[-1pt]{\footnotesize [-0.176, 0.112]}\end{tabular} & \begin{tabular}[c]{@{}c@{}}-0.110\\[-1pt]{\footnotesize [-0.270, 0.040]}\end{tabular} & \begin{tabular}[c]{@{}c@{}}-0.257\\[-1pt]{\footnotesize [-0.447, -0.071]}\end{tabular} & \begin{tabular}[c]{@{}c@{}}-0.139\\[-1pt]{\footnotesize [-0.311, 0.022]}\end{tabular} \\
이익을 낸 비율 & \begin{tabular}[c]{@{}c@{}}-0.043\\[-1pt]{\footnotesize [-0.182, 0.081]}\end{tabular} & \begin{tabular}[c]{@{}c@{}}0.044\\[-1pt]{\footnotesize [-0.084, 0.183]}\end{tabular} & \begin{tabular}[c]{@{}c@{}}0.012\\[-1pt]{\footnotesize [-0.125, 0.150]}\end{tabular} & \begin{tabular}[c]{@{}c@{}}0.074\\[-1pt]{\footnotesize [-0.066, 0.223]}\end{tabular} \\
손실이 아닌 비율 & \begin{tabular}[c]{@{}c@{}}-0.074\\[-1pt]{\footnotesize [-0.208, 0.052]}\end{tabular} & \begin{tabular}[c]{@{}c@{}}0.078\\[-1pt]{\footnotesize [-0.055, 0.217]}\end{tabular} & \begin{tabular}[c]{@{}c@{}}0.047\\[-1pt]{\footnotesize [-0.090, 0.191]}\end{tabular} & \begin{tabular}[c]{@{}c@{}}0.108\\[-1pt]{\footnotesize [-0.034, 0.256]}\end{tabular} \\
\bottomrule
\end{tabular}
}
\end{table}


% !TEX root = ../main.tex
\section{등록 분석}
\label{sec:map-registered}

등록 분석(미리 등록한 그대로 한 분석)은 분석 계획과 W1 등록 문서가 정해 둔 분석이에요. 이 연구가 판정에 쓸 자료를 뽑기 전에 정했어요(부록~\ref{sec:pipeline-stages}). 이 분석은 금액을 토큰의 원래 단위(소수점 조정을 하지 않은 가장 작은 단위)로 재고, 모든 ERC-20(토큰 공통 규칙) 토큰을 쓰고, 값 변환(금액이 아무리 커도 1을 넘지 않게 바꾸는 식)에 $b=1$을 써요. 또 2023년 10월 1일부터의 모든 이벤트(계약이 남기는 기록)로 점수를 매기고, 계약 $27{,}844$개로 된 계획의 코호트(묶음)를 써요. 수정 분석(등록 뒤에 금액을 달러로 바꿔 다시 한 분석)은 금액을 달러(USD)로 매기고, 다섯 토큰을 쓰고, 동결일(점수를 고정하는 날)마다 그 앞 열두 달의 기록으로 점수를 매겨요. 수정 분석은 등록 분석의 결과를 안 뒤에 정했어요(\ref{sec:amount-revision}절). 그리고 본문의 장들은 수정 분석을 보고해요. 표~\ref{tab:registered-vs-revised}는 등록 분석의 판정 대비와 판정을, 고친 점수에서 구한 같은 대비 옆에 나란히 놓아요. 이 표의 마지막 네 행은 라벨(나중에 채점에 쓰는 정답 값)을 본 뒤에 계산한 차수 대비이고, 판정이 없어요. 두 열은 R1을 빼면 같은 판정에 이르러요. R1은 고친 점수에서는 성립하고, 등록 분석에서는 성립하지 않아요. 이 연구가 판정에 쓸 자료를 뽑기 전에 정한 시험의 결과는 등록 열의 판정뿐이에요. 모든 계수, 구간, 판정을 담은 등록 분석의 전체 결과는 바꾸지 않은 그대로 \texttt{data/3-published-results-for-thesis/}에 공개했어요. 그 요약 파일인 \texttt{results\_digest.md}도 함께 공개했어요.

% Generated by experiments/scripts/export_latex.py -- do not edit by hand.
% Source: experiments/data/2-processed-tables-and-evaluations/same-window/eval_summary.json
% Source: experiments/data/2-processed-tables-and-evaluations/holdout-w0/eval_summary.json
% Source: experiments/data/2-processed-tables-and-evaluations/registered-w1/eval_summary.json
% Source: experiments/data/2-processed-tables-and-evaluations/revised-usd-12m/same-window/eval_summary.json
% Source: experiments/data/2-processed-tables-and-evaluations/revised-usd-12m/holdout-w0/eval_summary.json
% Source: experiments/data/2-processed-tables-and-evaluations/revised-usd-12m/window-w1/eval_summary.json
\begin{table}[htbp]
\centering
\caption[등록 분석과 수정 분석의 판정 대비.]{등록 분석의 판정 대비와 판정이에요(금액은 원래 단위, 모든 ERC-20 토큰, $b=1$, 점수는 2023년 10월 1일부터의 모든 이벤트로 매김). 분석 계획과 W1 등록 문서에 따라 얻은 그대로예요. 그 옆에 고친 점수에서 구한 같은 대비를 놓았어요(금액은 달러(USD), 목록에 오른 다섯 토큰, 점수는 각 동결일 전 열두 달의 이벤트로 매김). 수정 분석은 점수 기간까지 포함해서, 등록 분석의 결과를 안 뒤에 정했어요. 수정 분석 열은 등록한 규칙과 기준값을 고친 점수에 적용한 것이고, 등록한 검사가 아니에요. 마지막 묶음에는 두 분석의 어떤 규칙도 판정에 쓰지 않는 차수 대비가 있어요. 이 대비들은 두 분석 모두에서 보고해요. $\Delta\tau$와 95\% 짝지은 부트스트랩 구간을 적었어요(다시 뽑기 400번과 2{,}000번, 시드 42). 규칙과 그 기준값은 \ref{sec:holdout-protocol}절과 \ref{sub:fresh-holdout}절에 있어요.}
\label{tab:registered-vs-revised}
\fitwidth{%
\begin{tabular}{lcc}
\toprule
항목 & 등록 분석 & 수정 분석 \\
\midrule
매칭 코호트 (계약 수) & 27{,}844 & 6{,}483 \\
W0 스펜더 코호트 & 33{,}101 & 7{,}914 \\
W1 스펜더 코호트 & 34{,}587 & 7{,}960 \\
W1 트레이더 & 70 & 70 \\
EndorseRank 대 AWP, 같은 기간 ($\tau$) & 0.217 [0.209, 0.225] & 0.199 [0.181, 0.221] \\
\midrule
A, W0: C-PR 빼기 C-PR ($\lambda=1$), 새 허락 쌍 & $+0.008$ [0.003, 0.013] & $+0.031$ [0.021, 0.042] \\
A, W0: C-PR 빼기 C-PR ($\lambda=0$), 새 송금자 & $-0.030$ [-0.034, -0.026] & $-0.030$ [-0.040, -0.020] \\
규칙 A, 주된 부분 & 충족 안 됨 & 충족 안 됨 \\
규칙 A, 보조 부분 & 충족 안 됨 & 충족 안 됨 \\
\midrule
B, F1: C-PR 빼기 C-PR ($\lambda=0$), 새 허락 쌍 & $+0.101$ [0.092, 0.111] & $+0.170$ [0.149, 0.189] \\
B, F2: C-PR 빼기 C-PR ($\lambda=0$), 새 송금자 & $-0.025$ [-0.029, -0.021] & $-0.025$ [-0.038, -0.015] \\
B, F3: C-PR 빼기 C-PR ($\lambda=0$), 청산 없이 닫은 비율 & $-0.044$ [-0.100, -0.002] & $+0.000$ [-0.010, 0.010] \\
규칙 B (F1, F2, F3) & 충족 안 됨 & 충족 안 됨 \\
규칙 B, 민감도 1 (AWP를 비교 대상으로 함) & 충족 안 됨 & 충족 안 됨 \\
규칙 B, 민감도 2 (외톨이 점) & 충족 안 됨 & 충족 안 됨 \\
F6 (규칙 A, 주된 부분, W1에서) & 불성립 & 불성립 \\
\midrule
P, W1: 받은 허락 수 빼기 송금 입차수 & $+0.063$ [-0.026, 0.148] & $+0.141$ [0.037, 0.240] \\
Q, W1: EndorseRank (활동 재시작) 빼기 AWP & $+0.051$ [-0.085, 0.213] & $+0.023$ [-0.062, 0.113] \\
R1 (그냥 센 수의 수준에서 앞섬) & 충족 안 됨 & 충족됨 \\
R2 (페이지랭크 수준의 앞섬) & 충족 안 됨 & 충족 안 됨 \\
R3 (특이성) & 적용됨 & 적용됨 \\
\midrule
W0: EndorseRank 빼기 받은 허락 수, 새 쌍 & $-0.011$ [-0.016, -0.006] & $-0.090$ [-0.100, -0.079] \\
W0: AWP 빼기 송금 입차수, 새 송금자 & $+0.038$ [0.035, 0.042] & $-0.038$ [-0.043, -0.032] \\
W1: EndorseRank 빼기 받은 허락 수, 새 쌍 & $-0.029$ [-0.032, -0.025] & $-0.082$ [-0.090, -0.073] \\
W1: AWP 빼기 송금 입차수, 새 송금자 & $+0.020$ [0.016, 0.023] & $-0.041$ [-0.047, -0.034] \\
\bottomrule
\end{tabular}
}
\end{table}


% !TEX root = ../main.tex
\section{윤리 승인 증명서}
\label{ch:appendix-ethics}

\dissertationSubheading[sec:ethics-certificate]{증명서}

이 연구는 누구나 볼 수 있는 공개 블록체인 자료만 써요. 그리고 사람을 연구 참여자로 두지 않아요(\ref{sec:ethics}절). 아래에는 유니사(Unisa) 컴퓨팅 학부 / 농업·환경 과학 대학의 연구 윤리 심사 위원회가 발급한 윤리 승인(연구가 윤리 기준을 지킨다는 허가) 증명서를 옮겨 실어요.

% TODO (user action): replace the placeholder box with the certificate PDF once
% the reference number, date and scanned document are available, e.g.
%   \includegraphics[width=\textwidth,page=1]{Figures/ethics-clearance.pdf}
% Do not invent a reference number or date.
\begin{center}
\fbox{\parbox{0.9\textwidth}{\centering\vspace{2em}
윤리 승인 증명서\\[0.5em]
참조 번호: [Unisa SoC/CAES 윤리 승인 참조 번호: TBD]\\
발급일: [TBD]\\[0.5em]
(스캔한 증명서를 여기에 넣을 예정이에요.)
\vspace{2em}}}
\end{center}

